\SGAchapter{XX}{半单秩为 1 的约化群}{Groupes réductifs de rang semi-simple 1}
\begin{center}M. Demazure\end{center}
\SourcePage{27}
\SGAsection{1}{初等系统；群 $U_\alpha$ 与 $U_{-\alpha}$}
\textbf{回顾 1.1。}设 $S=\operatorname{Spec}(k)$，其中 $k$ 为\Term{algebraically-closed-field}{代数闭域}，$G$ 为\Term{semisimple-rank}{半单秩}为 $1$ 的 $S$-\Term{reductive-group}{约化群}，$T$ 为 $G$ 的一个\Term{maximal-torus}{极大环面}（必为\Term{split-torus}{分裂环面}）。则
\[\mathfrak g=\mathfrak t\oplus\mathfrak g^\alpha\oplus\mathfrak g^{-\alpha},\]
其中 $\alpha,-\alpha$ 为 $G$ 相对于 $T$ 的\Term{root}{根}。此外存在两个群\Term{monomorphism}{单态射}
\[p_\alpha:\mathbf G_{a,S}\longrightarrow G,\qquad p_{-\alpha}:\mathbf G_{a,S}\longrightarrow G,\]
使得对任意 $S'\to S$ 以及 $t\in T(S')$、$x\in\mathbf G_a(S')$，有
\[t p_\alpha(x)t^{-1}=p(\alpha(t)x),\qquad t p_{-\alpha}(x)t^{-1}=p_{-\alpha}(\alpha(t)^{-1}x),\]
而由 $(y,t,x)\mapsto p_{-\alpha}(y)t p_\alpha(x)$ 定义的\Term{morphism}{态射}
\[\mathbf G_{a,S}\times_S T\times_S\mathbf G_{a,S}\longrightarrow G\]
为\Term{radicial}{根性}且\Term{dominant}{占优}（Bible，§13.4，定理3的推论2）。由于单位元处的切映射为\Term{bijection}{双射}，该\Term{morphism}{态射}也\Term{separable}{可分}，故\Term{birational}{双有理}；由 Zariski 的“主定理”（EGA III$_1$，4.4.9），它是\Term{open-immersion}{开浸入}。
\begin{SGAresult}{引理}{1.2}
设 $S$ 为\Term{scheme}{概形}，$G$ 为 $S$-\Term{group-scheme}{群概形}，$T$ 为 $G$ 的\Term{torus}{环面}，$Q$ 为 $T$ 的\Term{subtorus}{子环面}，$\alpha$ 为 $T$ 的\Term{character}{特征标}，在 $Q$ 的每个\Term{fiber}{纤维}上均诱导非平凡的\Term{character}{特征标}。设 $p_\alpha:\mathbf G_{a,S}\to G$（相应地 $p_{-\alpha}$）为受 $T$ \Term{normalization}{正规化}、\Term{multiplier}{乘子}为 $\alpha$（相应地 $-\alpha$）的\Term{group-morphism}{群态射}。假定由 $u(y,t,x)=p_{-\alpha}(y)t p_\alpha(x)$ 定义的\Term{morphism}{态射}
\[u:\mathbf G_{a,S}\times_S T\times_S\mathbf G_{a,S}\longrightarrow G\]
为\Term{open-immersion}{开浸入}。再设 $q\geq0$ 为整数，$p:\mathbf G_{a,S}\to G$ 为\Term{group-morphism}{群态射}，
\SourcePage{28}
使得对任意 $S'\to S$、$t\in Q(S')$、$x\in\mathbf G_a(S')$，有
\[\operatorname{int}(t)^q(p(x))=p(\alpha(t)x).\]
则存在唯一 $\nu\in\mathbf G_a(S)$，使 $p(x)=p_\alpha(\nu x^q)$。
\end{SGAresult}
\begin{proof}
设 $\Omega$ 为 $u$ 的像，$U=p^{-1}(\Omega)$。这是 $\mathbf G_{a,S}$ 中包含\Term{zero-section}{零截面}的\Term{open-subset}{开集}。对 $Q$ 的任意\Term{section}{截面} $t$，以 $\alpha(t)$ 相乘定义的 $\mathbf G_{a,S}$ \Term{automorphism}{自同构}整体保持 $U$。有 $U=\mathbf G_{a,S}$：只须在 $S$ 为\Term{algebraically-closed-field}{代数闭域} $k$ 的谱时验证；此时 $\alpha:Q(k)\to k^*$ \Term{surjectivity}{满射}，立即给出 $U(k)\supset k^*$，故 $U=\mathbf G_{a,k}$。因此存在三个\Term{morphism}{态射}
\[a:\mathbf G_{a,S}\to\mathbf G_{a,S},\quad b:\mathbf G_{a,S}\to T,\quad c:\mathbf G_{a,S}\to\mathbf G_{a,S}\]
使 $p(x)=p_{-\alpha}(a(x))b(x)p_\alpha(c(x))$。关于 $p$ 的条件成为
\[a(\alpha(t)x)=\alpha(t)^{-q}a(x),\quad b(\alpha(t)x)=b(x),\quad c(\alpha(t)x)=\alpha(t)^q c(x).\]
由同一理由，对任意 $S'\to S$ 和 $z\in\mathbf G_m(S')$，有
\[a(zx)=z^{-q}a(x),\quad b(zx)=b(x),\quad c(zx)=z^q c(x),\]
从而 $z^q a(z)=a(1)$、$b(z)=b(1)$、$c(z)=z^q c(1)$。由于 $\mathbf G_{m,S}$ 在 $\mathbf G_{a,S}$ 中\Term{schematically-dense}{概形稠密}，对所有 $x\in\mathbf G_a(S')$ 有
\[\begin{aligned}x^q a(x)&=a(1)=a(0)=0,&\text{故 }a&=0,\\c(x)&=x^q c(1)=\nu x^q,&\nu&\in\mathbf G_a(S),\\b(x)&=b(1)=b(0)=e,&\text{故 }b&=e.\end{aligned}\]
证明完毕。
\end{proof}
\begin{SGAresult}{定义}{1.3}
设 $S$ 为\Term{scheme}{概形}。称三元组 $(G,T,\alpha)$ 为 $S$-\Term{elementary-system}{初等系统}，若：(i) $G$ 为\Term{semisimple-rank}{半单秩}为 $1$ 的 $S$-\Term{reductive-group}{约化群}（第 XIX 讲2.7）；(ii) $T$ 为 $G$ 的\Term{maximal-torus}{极大环面}；(iii) $\alpha$ 为 $G$ 相对于 $T$ 的\Term{root}{根}（第 XIX 讲3.2）。
\end{SGAresult}
于是有\Term{direct-sum}{直和}\Term{resolution}{分解}（第 XIX 讲3.5）\footnote{编者注：这是 $\mathscr O_S$-模的\Term{direct-sum}{直和}。}
\[\mathfrak g=\mathfrak t\oplus\mathfrak g^\alpha\oplus\mathfrak g^{-\alpha},\]
其中 $\mathfrak g^\alpha,\mathfrak g^{-\alpha}$ 均\Term{locally-free}{局部自由}且秩为一。
\SourcePage{29}
\textbf{1.4。}若 $(G,T,\alpha)$ 为 $S$-\Term{elementary-system}{初等系统}，则对任意 $S'\to S$，$(G_{S'},T_{S'},\alpha_{S'})$ 为 $S'$-\Term{elementary-system}{初等系统}；$(G,T,-\alpha)$ 也是 $S$-\Term{elementary-system}{初等系统}。

若 $S$ 为\Term{scheme}{概形}，$G$ 为 $S$-\Term{reductive-group}{约化群}，$T$ 为 $G$ 的\Term{maximal-torus}{极大环面}，$\alpha$ 为 $G$ 相对于 $T$ 的\Term{root}{根}，则由第 XIX 讲3.9，$(Z_\alpha,T,\alpha)$ 是 $S$-\Term{elementary-system}{初等系统}。

设 $(G,T,\alpha)$ 为 $S$-\Term{elementary-system}{初等系统}。\Term{invertible-module}{可逆模} $\mathfrak g^\alpha$ 具有\Term{canonical}{典范}的 $T$-模结构，故\Term{vector-bundle}{向量丛} $\mathbf W(\mathfrak g^\alpha)$ 也具有 $T$-模结构。另一方面，$T$ 的\Term{inner-automorphism}{内自同构}在 $G$ 上定义以 $T$ 为\Term{operator}{算子}群的\Term{group-structure}{群结构}。
\begin{SGAresult}{定理}{1.5}
(i) 存在唯一的以 $T$ 为\Term{operator}{算子}群的\Term{group-morphism}{群态射}
\[\exp:\mathbf W(\mathfrak g^\alpha)\longrightarrow G,\]
在\Term{lie-algebra}{李代数}上诱导\Term{canonical}{典范}\Term{morphism}{态射} $\mathfrak g^\alpha\to\mathfrak g$。\footnote{编者注：下文推论5.9将说明，$\exp$ 是 $\mathbf W(\mathfrak g^\alpha)$ 到 $G$ 的一个\Term{closed-subgroup}{闭子群}的\Term{isomorphism}{同构}。}换言之，$\exp$ 是唯一满足下列条件的\Term{morphism}{态射}：对任意 $S'\to S$、$t\in T(S')$ 及 $X,X'\in\mathbf W(\mathfrak g^\alpha)(S')$，
\[\exp(X+X')=\exp(X)\exp(X'),\quad \operatorname{int}(t)(\exp(X))=\exp(\alpha(t)X),\quad\operatorname{Lie}(\exp)(X)=X.\]
(ii) 同样在 $S$-\Term{elementary-system}{初等系统} $(G,T,-\alpha)$ 中定义 $\exp:\mathbf W(\mathfrak g^{-\alpha})\to G$，则由 $(Y,t,X)\mapsto\exp(Y)t\exp(X)$ 定义的\Term{morphism}{态射}
\[\mathbf W(\mathfrak g^{-\alpha})\times_S T\times_S\mathbf W(\mathfrak g^\alpha)\longrightarrow G\]
是\Term{open-immersion}{开浸入}。
\end{SGAresult}
\begin{proof}
先假设所求 $\exp$ 已存在，证明其余断言。先证(ii)。两端均在 $S$ 上\Term{finitely-presented}{有限呈示}且\Term{flat}{平坦}，故只须在 $S$ 为\Term{algebraically-closed-field}{代数闭域}的谱时证明（SGA 1，I 5.7及 VIII 5.5）。设 $S=\operatorname{Spec}k$，取 $X\in\Gamma(S,\mathfrak g^\alpha)^\times$、$Y\in\Gamma(S,\mathfrak g^{-\alpha})^\times$，只须证明 $(y,t,x)\mapsto\exp(yY)t\exp(xX)$ 是\Term{open-immersion}{开浸入}。由1.1及1.2，存在 $a,b\in k$，使
\[\exp(yY)=p_{-\alpha}(ay),\qquad\exp(xX)=p_\alpha(bx).\]
由于 $\exp:\mathbf W(\mathfrak g^{-\alpha})\to G$ 在\Term{lie-algebra}{李代数}上诱导\Term{monomorphism}{单态射}，$a\ne0$；同样 $b\ne0$，遂归结于1.1。
\SourcePage{30}
$\exp$ 的唯一性可在 $S$ 上\Term{local-scheme}{局部}证明，遂可假设 $\mathfrak g^\alpha,\mathfrak g^{-\alpha}$ 自由，再用1.2（$Q=T,q=1$）。

只须证明存在性。由\Term{faithfully-flat}{忠实平坦}\Term{descent}{下降}及上述唯一性，只须对(fpqc)\Term{topology}{拓扑}在 $S$ 上\Term{local-scheme}{局部}证明。用通常的\Term{finitely-presented}{有限呈示}论证，先归结于 $S$ \Term{noetherian}{诺特}，再归结于\Term{noetherian}{诺特}\Term{local-scheme}{局部}情形。因此只须在 $S=\operatorname{Spec}(A)$、$A$ 为\Term{residue-field}{剩余域} $k$ \Term{algebraically-closed}{代数闭}的\Term{complete}{完备}\Term{noetherian-local-ring}{诺特局部环}时证明。取受 $T_k$ \Term{normalization}{正规化}、\Term{multiplier}{乘子}为 $\alpha_0=\alpha\otimes_Ak$ 的 $k$-群\Term{monomorphism}{单态射} $p_0:\mathbf G_{a,k}\to G_k$（由1.1存在）。由1.1及1.2，相应的 $T_k\mathbin{\cdot}_{\alpha_0}\mathbf G_{a,k}\to G_k$ 为\Term{immersion}{浸入}，尤其为\Term{monomorphism}{单态射}。暂且承认以下两引理。
\end{proof}
\begin{SGAresult}{引理}{1.6}
设 $S$ 为\Term{scheme}{概形}，$G$ 为\Term{finitely-presented}{有限呈示}的 $S$-群，$T$ 为 $S$-\Term{torus}{环面}，$\alpha$ 为 $T$ 的一个在每个\Term{fiber}{纤维}上非平凡的\Term{character}{特征标}，$s_0\in S$。设 $f:T\mathbin{\cdot}_\alpha\mathbf G_{a,S}\to G$ 为 $S$-\Term{group-morphism}{群态射}，$f_{s_0}$ 为\Term{monomorphism}{单态射}，且 $f$ 在 $T$ 上的\Term{restriction}{限制}也是\Term{monomorphism}{单态射}。则存在 $s_0$ 的\Term{open-neighborhood}{开邻域} $U$，使 $f|_U$ 为\Term{monomorphism}{单态射}。
\end{SGAresult}
\begin{SGAresult}{引理}{1.7}
设 $A$ 为\Term{residue-field}{剩余域} $k$ \Term{algebraically-closed}{代数闭}的\Term{complete}{完备}\Term{noetherian-local-ring}{诺特局部环}，$(G,T,\alpha)$ 为 $A$-\Term{elementary-system}{初等系统}，$p_0:\mathbf G_{a,k}\to G_k$ 为受 $T_k$ \Term{normalization}{正规化}、\Term{multiplier}{乘子}为 $\alpha\otimes_Ak$ 的 $k$-\Term{group-morphism}{群态射}。则存在受 $T$ \Term{normalization}{正规化}、\Term{multiplier}{乘子}为 $\alpha$ 的\Term{group-morphism}{群态射} $p:\mathbf G_{a,A}\to G$。
\end{SGAresult}
取1.7所给的 $p$，以及相应的 $f:T\mathbin{\cdot}_\alpha\mathbf G_{a,S}\to G$。它满足1.6的假设，故为\Term{monomorphism}{单态射}；尤其 $p$ 为\Term{monomorphism}{单态射}。于是由第 XIX 讲4.9得结论。
\begin{proof}[1.6的证明]
以 $\varepsilon:S\to T\mathbin{\cdot}_\alpha\mathbf G_{a,S}$ 为\Term{unit-section}{单位截面}。$f$ 在 $\varepsilon(s_0)$ 处\Term{unramified}{非分歧}，故对 $s_0$ 某\Term{open-neighborhood}{开邻域} $U$ 内所有 $s$，在 $\varepsilon(s)$ 处均\Term{unramified}{非分歧}；由第 X 讲3.5，$f|_U$ \Term{unramified}{非分歧}，故核 $\operatorname{Ker}(f)_U$ 在 $U$ 上\Term{unramified}{非分歧}。因此要证明 $f|_U$ 为\Term{monomorphism}{单态射}，只须\footnote{编者注：据 EGA IV$_4$，17.9.1。}证明该核在 $U$ 上\Term{radicial}{根性}，这是\Term{set-theory}{集合论}问题。遂归结于下述引理。
\end{proof}
\begin{SGAresult}{引理}{1.8}
设 $k$ \Term{algebraically-closed}{代数闭}，$N$ 为 $T\mathbin{\cdot}_\alpha\mathbf G_{a,k}$ 的\Term{normal-subgroup}{正规子群}（$\alpha$ 为\Term{torus}{环面} $T$ 的非平凡\Term{character}{特征标}），在 $k$ 上\Term{etale}{平展}且 $N\cap T=\{e\}$。则 $N=\{e\}$。
\end{SGAresult}
\begin{proof}
有 $\operatorname{int}(t')(t,x)=(t,\alpha(t')x)$。若 $(t,x)\in N$ 且 $x\ne0$，则对 $z\in k^*$，$(t,zx)$ 也在 $N$ 中；于是 $(t,x)$ 非孤立，$N$ 非\Term{quasi-finite}{拟有限}。故\Term{set}{集合}意义下 $N\subset T$，得证。
\end{proof}
\SourcePage{31}
\begin{proof}[1.7的证明]
令 $\mathfrak m$ 为 $A$ 的\Term{jacobson-radical}{根基}，$S_n=\operatorname{Spec}(A/\mathfrak m^{n+1})$，$n\geq0$。先对 $n$ 归纳，证明 $p_0$ 可延拓为 $S_n$-\Term{group-morphism}{群态射} $p_n:\mathbf G_{a,S_n}\to G_{S_n}$，受 $T_{S_n}$ \Term{normalization}{正规化}，\Term{multiplier}{乘子}为 $\alpha_n$，并且 $p_{n+1}\times_{S_{n+1}}S_n=p_n$。

令 $H=T\mathbin{\cdot}_\alpha\mathbf G_{a,S}$，$f_n$ 为 $p_n$ 所定义的 $H_{S_n}\to G_{S_n}$。暂且承认以下引理。
\end{proof}
\begin{SGAresult}{引理}{1.9}
若 $(G,T,\alpha)$ 是 $k$-\Term{elementary-system}{初等系统}，$k$ \Term{algebraically-closed}{代数闭}，$p:\mathbf G_{a,k}\to G$ 为受 $T$ \Term{normalization}{正规化}、\Term{multiplier}{乘子}为 $\alpha$ 的\Term{monomorphism}{单态射}，则
\[\mathrm H^2(T\mathbin{\cdot}_\alpha\mathbf G_{a,k},\mathfrak g)=0.\]
这里 $T\mathbin{\cdot}_\alpha\mathbf G_{a,k}$ 经由 $p$ 所定义的到 $G$ 的\Term{morphism}{态射}及 $G$ 的\Term{adjoint-representation}{伴随表示}\Term{action}{作用}于 $\mathfrak g$。
\end{SGAresult}
由第 III 讲2.8，$f_n$ 延拓为 $S_{n+1}$-\Term{group-morphism}{群态射} $f'_{n+1}:H_{S_{n+1}}\to G_{S_{n+1}}$。$f'_{n+1}$ 与 $T_{S_{n+1}}$ 到 $G_{S_{n+1}}$ 的\Term{canonical}{典范}\Term{immersion}{浸入}在 $T_{S_n}$ 上\Term{restriction}{限制}相同。由第 III 讲2.5，存在 $g\in G(S_{n+1})$，$g\times_{S_{n+1}}S_n=e$，使 $f_{n+1}=\operatorname{int}(g)\circ f'_{n+1}$ 在 $T_{n+1}$ 上\Term{restriction}{限制}为\Term{canonical}{典范}\Term{immersion}{浸入}。以 $p_{n+1}$ 为 $f_{n+1}$ 在 $\mathbf G_{a,S_{n+1}}$ 上的\Term{restriction}{限制}；它受 $T_{S_{n+1}}$ \Term{normalization}{正规化}，\Term{multiplier}{乘子}为 $\alpha_{S_{n+1}}$，并延拓 $p_n$。

这样构造了一相容系统 $(f_n)$；现须将其\Term{algebraization}{代数化}。我们有：
\begin{SGAresult}{引理}{1.10}
设 $A$ 为\Term{complete}{完备}\Term{noetherian-local-ring}{诺特局部环}，$\mathfrak m$ 为\Term{maximal-ideal}{极大理想}，$S=\operatorname{Spec}(A)$，$S_n=\operatorname{Spec}(A/\mathfrak m^{n+1})$，$T$ 为 $S$-\Term{torus}{环面}，$\alpha$ 为其非零\Term{character}{特征标}，$X$ 为带 $T$-\Term{action}{作用}的\Term{affine}{仿射} $S$-\Term{scheme}{概形}。令 $T$ 经 $\alpha$ \Term{action}{作用}于 $\mathbf G_{a,S}$。设 $q\geq0$，$(f_n)_{n\geq0}$ 为以 $T_{S_n}$ 为\Term{operator}{算子}的\Term{object}{对象}\Term{morphism}{态射} $f_n:\mathbf G_{a,S_n}^q\to X_{S_n}$ 的相容系统。则存在唯一以 $T$ 为\Term{operator}{算子}的\Term{object}{对象}\Term{morphism}{态射} $f:\mathbf G_{a,S}^q\to X$ 诱导诸 $f_n$（比较第 IX 讲7.1）。
\end{SGAresult}
\begin{SGAresult}{推论}{1.11}
若 $X$ 是以 $T$ 为\Term{operator}{算子}群的群，诸 $f_n$ 为\Term{group-morphism}{群态射}，则 $f$ 亦为\Term{group-morphism}{群态射}。
\end{SGAresult}
\begin{proof}只须把引理的唯一性断言用于通常由 $f$ \Term{derive}{导出}的两个\Term{morphism}{态射} $\mathbf G_{a,S}^{2q}\to X$。\end{proof}
\begin{proof}[1.10的证明]
设 $T$ \Term{splitting}{分裂}；这也正是在1.5的证明中应用1.10时的情形。由第 I 讲4.7.3的注，$X\mapsto\mathscr A(X)$ 给出带 $T$-\Term{action}{作用}的\Term{affine}{仿射} $S$-\Term{scheme}{概形}的\Term{category}{范畴}与 $M=\operatorname{Hom}_{S\text{-gr.}}(T,\mathbf G_{m,S})$ 型\Term{graded}{分次} $S$-\Term{algebra}{代数}\Term{category}{范畴}之\Term{opposite-category}{反范畴}的\Term{equivalence}{等价}。
\SourcePage{32}
故有\Term{graded}{分次}
\[B=\mathscr A(X)=\coprod_{m\in M}B_m,\qquad C=\mathscr A(\mathbf G_{a,S}^q)=\coprod_{m\in M}C_m.\]
每个 $C_m$ 显然在 $A$ 上自由且\Term{finite-type}{有限型}。（若 $m$ 非 $\alpha$ 的倍数则 $C_m=0$；若 $m=d\alpha$，则 $C_m$ \Term{isomorphism}{同构}于 $q$ 变元 $d$ 次\Term{homogeneous}{齐次}多项式的 $A$-模。）令
\[\widehat B_m=\varprojlim_n B_m\otimes_A(A/\mathfrak m^{n+1}),\quad\widehat C_m=\varprojlim_n C_m\otimes_A(A/\mathfrak m^{n+1}),\quad\widehat B=\coprod_{m\in M}\widehat B_m,\quad\widehat C=\coprod_{m\in M}\widehat C_m.\]
有\Term{canonical}{典范}的 $M$ 型\Term{graded-algebra}{分次代数}\Term{morphism}{态射} $g_B:B\to\widehat B$、$g_C:C\to\widehat C$；由以上观察，$g_C$ 为\Term{isomorphism}{同构}。给定引理中的相容系统 $(f_n)$ \Term{equivalence}{等价}于给定\Term{graded}{分次} $A$-\Term{algebra}{代数}\Term{morphism}{态射} $\widehat F:\widehat B\to\widehat C$。寻找所求 $f$ \Term{equivalence}{等价}于寻找使下图交换的\Term{graded}{分次} $A$-\Term{algebra}{代数}\Term{morphism}{态射} $F:B\to C$：
\[\begin{tikzcd}B\arrow[r,"F"]\arrow[d,"g_B"']&C\arrow[d,"g_C"]\\\widehat B\arrow[r,"\widehat F"']&\widehat C.\end{tikzcd}\]
$g_C$ 为\Term{isomorphism}{同构}，故 $F$ 的存在性与唯一性显然。1.10得证。
\end{proof}
为完成1.5的证明，尚须证明1.9。

\textbf{1.12。1.9的证明。}有 $\mathfrak g=\mathfrak t\oplus\mathfrak g^\alpha\oplus\mathfrak g^{-\alpha}$。按1.9将 $\mathfrak g$ 视为 $(T\mathbin{\cdot}_\alpha\mathbf G_{a,k})$-模。显然 $\mathfrak t\oplus\mathfrak g^\alpha$ 为\Term{submodule}{子模}，商作为 $k$-\Term{vector-space}{向量空间}，甚至作为 $T$-模，\Term{isomorphism}{同构}于 $\mathfrak g^{-\alpha}$。$\mathbf G_{a,k}$ 在此一维商上的\Term{action}{作用}显然平凡（$\mathbf G_{a,k}\to\mathbf G_{m,k}$ 的任何\Term{group-morphism}{群态射}均平凡）。同样 $\mathfrak g^\alpha$ 为 $\mathfrak t\oplus\mathfrak g^\alpha$ 的\Term{submodule}{子模}，商作为 $T$-模\Term{isomorphism}{同构}于 $\mathfrak t$，$\mathbf G_{a,k}$ 在其上\Term{action}{作用}平凡。综上：
\begin{SGAresult}{引理}{1.13}
在1.9的条件下，$\mathfrak g$ 作为 $(T\mathbin{\cdot}_\alpha\mathbf G_{a,k})$-模有一组成列，相继商为 $\mathfrak g^{-\alpha},\mathfrak t,\mathfrak g^\alpha$，均经\Term{projection}{投影} $T\mathbin{\cdot}_\alpha\mathbf G_{a,k}\to T$ 视为模。
\end{SGAresult}
\SourcePage{33}
因此归结为计算 $T\mathbin{\cdot}_\alpha\mathbf G_{a,k}$ 经\Term{projection}{投影}到 $T$ 及 $T$ 的\Term{character}{特征标} $\beta$（此处 $\beta=0,\alpha,-\alpha$）\Term{action}{作用}于 $\mathbf W(k)$ 的\Term{cohomology}{上同调}。\footnote{编者注：添加了下一句。}记 $k[x_1,\ldots,x_n]$ 为 $n$ 变元多项式\Term{algebra}{代数}，$k_q[x_1,\ldots,x_n]$ 为 $q$ 次\Term{homogeneous}{齐次}多项式子\Term{space}{空间}。
\begin{SGAresult}{引理}{1.14}
在以上记号下，$\mathrm H^n(T\mathbin{\cdot}_\alpha\mathbf G_{a,k},k)=\mathrm H^n(C^*_{\alpha,\beta})$，其中
\[C^n_{\alpha,\beta}=\begin{cases}k_q[x_1,\ldots,x_n]&\beta=q\alpha,\ q\in\mathbf N^*,\\0&\text{否则},\end{cases}\]
且
\[\begin{split}\delta f(x_1,\ldots,x_{n+1})={}&f(x_2,\ldots,x_{n+1})+\sum_{i=1}^n(-1)^if(x_1,\ldots,x_ix_{i+1},\ldots,x_{n+1})\\&+(-1)^{n+1}f(x_1,\ldots,x_n).\end{split}\]
\end{SGAresult}
\begin{proof}
由第 I 讲5.3.2，\Term{functor}{函子} $M\mapsto\mathrm H^0(T,M)$ 在 $T$-\Term{module-category}{模范畴}上\Term{exact}{正合}（$\mathrm H^q(T,-)$ 均为零）。故与通\Term{constant-group}{常群}\Term{cohomology}{上同调}一样，所求 $\mathrm H^n$ 可由 $\mathbf G_{a,k}$ 到 $k$ 的 $T$-不变\Term{cochain-complex}{上链复形}的第 $n$ \Term{cohomology}{上同调}计算，即\Term{cochain}{上链}满足
\[f(\alpha(t)x_1,\ldots,\alpha(t)x_n)=\beta(t)f(x_1,\ldots,x_n).\]
这正给出所述\Term{complex}{复形}。为证明1.9，只须证明 $\beta=0,\alpha,-\alpha$ 时 $\mathrm H^2(C^*_{\alpha,\beta})=0$，这是直接的。
\end{proof}
\textbf{注记 1.15。}$\beta=q\alpha$ 时可显式计算 $\mathrm H^n(C^*_{\alpha,\beta})$（见 M. Lazard，\textit{Lois de groupes et analyseurs}，Annales E.N.S.，1955）。特别地，$n>q$ 时 $\mathrm H^n(C^*_{\alpha,q\alpha})=0$。

\textbf{记号 1.16。}将\Term{canonical}{典范}\Term{immersion}{浸入} $\mathbf W(\mathfrak g^{-\alpha})\times_S T\times_S\mathbf W(\mathfrak g^\alpha)\to G$ 的像记为 $\Omega$。它是包含\Term{unit-section}{单位截面}的\Term{open-subset}{开集}。$\mathbf W(\mathfrak g^{-\alpha})$、$\mathbf W(\mathfrak g^\alpha)$、$\mathbf W(\mathfrak g^{-\alpha})\times_ST$、$T\times_S\mathbf W(\mathfrak g^\alpha)$ 的像分别记为\footnote{编者注：将 $P_{-\alpha},P_\alpha$ 改为 $U_{-\alpha},U_\alpha$。}
\[U_{-\alpha},\quad U_\alpha,\quad U_{-\alpha}\cdot T,\quad T\cdot U_\alpha.\]
$U_\alpha$（相应地 $U_{-\alpha}$）为 $G$ 的\Term{subgroup}{子群}，\Term{canonical}{典范}地带\Term{vector-bundle}{向量丛}结构，且对所有 $S'\to S$、$t\in T(S')$、$x\in U_\alpha(S')$（相应地 $x\in U_{-\alpha}(S')$），
\[\operatorname{int}(t)(x)=x^{\alpha(t)}\quad\text{（相应地 }x^{-\alpha(t)}\text{）。}\]
\SourcePage{34}
有\Term{canonical-isomorphism}{典范同构} $T\cdot U_\alpha\simeq T\mathbin{\cdot}_\alpha U_\alpha$、$T\cdot U_{-\alpha}\simeq T\mathbin{\cdot}_{-\alpha}U_{-\alpha}$。$\Omega$ 在 $\operatorname{int}(T)$ 下\Term{stable}{稳定}：
\[\operatorname{int}(t')(y\cdot t\cdot x)=y^{-\alpha(t')}\cdot t\cdot x^{\alpha(t')}.\]
\begin{SGAresult}{推论}{1.17}
$\operatorname{Lie}(U_\alpha/S)=\mathfrak g^\alpha$，$\operatorname{Lie}(U_{-\alpha}/S)=\mathfrak g^{-\alpha}$；\Term{isomorphism}{同构} $\exp:\mathbf W(\mathfrak g^\alpha)\xrightarrow\sim U_\alpha$ 及 $\exp:\mathbf W(\mathfrak g^{-\alpha})\xrightarrow\sim U_{-\alpha}$ 正是第 XIX 讲4.2中的\Term{isomorphism}{同构}。
\end{SGAresult}
\begin{SGAresult}{推论}{1.18}$\Omega$ 在 $G$ 中相对\Term{schematically-dense}{概形稠密}（参见第 XVIII 讲§1）。\end{SGAresult}
\begin{proof}由第 XVIII 讲1.3显然。\end{proof}
\begin{SGAresult}{推论}{1.19}
$G$ 的\Term{center}{中心}为 $\operatorname{Centr}(G)=\operatorname{Ker}(\alpha)$，故是 $G$ 的\Term{finite-type}{有限型}\Term{multiplicative-type}{乘法型}\Term{closed-subgroup}{闭子群}。
\end{SGAresult}
\begin{proof}
第二断言由第一断言及第 IX 讲2.7推出。$\operatorname{Ker}(\alpha)$ 的\Term{section}{截面}所定义的\Term{inner-automorphism}{内自同构}在 $\Omega$ 上\Term{action}{作用}平凡（1.16末式），由1.18在 $G$ 上也平凡。反之，若 $g\in G(S)$ \Term{centralizes}{中心化} $G$，则\Term{centralizes}{中心化} $T,U_\alpha$，故由第 XIX 讲2.8为 $T$ 的\Term{section}{截面}，且消去 $\alpha$。这也适用于任意\Term{base-change}{基变换}，故得所求等式。
\end{proof}
\begin{SGAresult}{推论}{1.20}
存在受 $T$ \Term{normalization}{正规化}、\Term{multiplier}{乘子}为 $\alpha$ 的\Term{monomorphism}{单态射} $p_\alpha:\mathbf G_{a,S}\to G$ 的充要条件是 $\mathscr O_S$-模 $\mathfrak g^\alpha$ 自由。更精确地，
\[X_\alpha\mapsto(x\mapsto\exp(xX_\alpha)),\qquad p_\alpha\mapsto\operatorname{Lie}(p_\alpha)\]
给出 $\Gamma(S,\mathfrak g^\alpha)^\times$ 与上述 $p_\alpha$ 的\Term{set}{集合}（也就是向量\Term{group-scheme}{群概形}\Term{isomorphism}{同构} $\mathbf G_{a,S}\xrightarrow\sim U_\alpha$ 的\Term{set}{集合}）之间的\Term{bijection}{双射}。\footnote{编者注：一方面 $\operatorname{Lie}(\mathbf G_{a,S})=\mathscr O_S$，另一方面 $\operatorname{Lie}(p_\alpha)\in\operatorname{Hom}_{\mathscr O_S}(\mathscr O_S,\mathfrak g^\alpha)=\Gamma(S,\mathfrak g^\alpha)$。}
\end{SGAresult}
\begin{SGAresult}{推论}{1.21}$G$ 的\Term{subgroup}{子群} $U_\alpha,U_{-\alpha}$ 在任何\Term{fiber}{纤维}上均不交换。\end{SGAresult}
\begin{proof}
否则若 $(U_\alpha)_s,(U_{-\alpha})_s$ 交换，$\Omega_s$ 就是 $G_s$ 的\Term{subgroup}{子群}，故 $\Omega_s=G_s$\footnote{编者注：据1.18。}且 $G_s$ \Term{solvable}{可解}，与其为\Term{semisimple-rank}{半单秩} $1$ 的\Term{reductive-group}{约化群}矛盾。
\end{proof}
\SourcePage{35}
\SGAsection{2}{初等系统的结构}
\begin{SGAresult}{定理}{2.1}
设 $S$ 为\Term{scheme}{概形}，$(G,T,\alpha)$ 为 $S$-\Term{elementary-system}{初等系统}。存在 $\mathscr O_S$-\Term{module-morphism}{模态射}
\[\mathfrak g^\alpha\otimes_{\mathscr O_S}\mathfrak g^{-\alpha}\longrightarrow\mathscr O_S,\quad(X,Y)\longmapsto\langle X,Y\rangle,\]
及 $S$-\Term{group-morphism}{群态射} $\alpha^*:\mathbf G_{m,S}\to T$，使对任意 $S'\to S$、$X\in\Gamma(S',\mathfrak g^\alpha\otimes\mathscr O_{S'})$、$Y\in\Gamma(S',\mathfrak g^{-\alpha}\otimes\mathscr O_{S'})$，有
\[\exp(X)\exp(Y)\in\Omega(S')\quad\Longleftrightarrow\quad1+\langle X,Y\rangle\in\mathbf G_m(S'),\]
而在此条件下
\[\exp(X)\exp(Y)=\exp\left(\frac Y{1+\langle X,Y\rangle}\right)\alpha^*(1+\langle X,Y\rangle)\exp\left(\frac X{1+\langle X,Y\rangle}\right).\tag{F}\]
此外，$(X,Y)\mapsto\langle X,Y\rangle$ 与 $\alpha^*$ 唯一确定，前者为\Term{isomorphism}{同构}，从而给出 $\mathfrak g^\alpha$ 与 $\mathfrak g^{-\alpha}$ 的\Term{duality}{对偶}，且 $\alpha\circ\alpha^*=2$（$\mathbf G_{m,S}$ 上的平方映射）。
\end{SGAresult}
由唯一性，只须在 $S$ 上\Term{local-scheme}{局部}证明。可设 $\mathfrak g^\alpha,\mathfrak g^{-\alpha}$ 自由。取 $X\in\Gamma(S,\mathfrak g^\alpha)^\times$、$Y\in\Gamma(S,\mathfrak g^{-\alpha})^\times$，对 $S'\to S$ 及 $x,y\in\mathbf G_a(S')$ 令 $p_\alpha(x)=\exp(xX)$、$p_{-\alpha}(y)=\exp(yY)$。由1.5及1.21，只须证明：
\begin{SGAresult}{命题}{2.2}
设 $S$ 为\Term{scheme}{概形}，$G$ 为 $S$-群，$T$ 为 $G$ 的\Term{torus}{环面}，$\alpha$ 为每个\Term{fiber}{纤维}上非平凡的 $T$ 的\Term{character}{特征标}，$p_\alpha:\mathbf G_{a,S}\to G$（相应地 $p_{-\alpha}$）为受 $T$ \Term{normalization}{正规化}、\Term{multiplier}{乘子}为 $\alpha$（相应地 $-\alpha$）的群\Term{monomorphism}{单态射}。假设：

(i) $(y,t,x)\mapsto p_{-\alpha}(y)t p_\alpha(x)$ 给出\Term{open-immersion}{开浸入} $\mathbf G_{a,S}\times_ST\times_S\mathbf G_{a,S}\to G$，像记为 $\Omega$；

(ii) 对每个 $s\in S$，$(p_\alpha)_s(\mathbf G_{a,\kappa(s)})$ 与 $(p_{-\alpha})_s(\mathbf G_{a,\kappa(s)})$ 不交换。

则存在唯一确定的 $a\in\mathbf G_a(S)$、$\alpha^*\in\operatorname{Hom}_{S\text{-gr.}}(\mathbf G_{m,S},T)$，使对任意 $S'\to S$、$x,y\in\mathbf G_a(S')$，
\[p_\alpha(x)p_{-\alpha}(y)\in\Omega(S')\quad\Longleftrightarrow\quad1+axy\in\mathbf G_m(S'),\]
且在此条件下
\[p_\alpha(x)p_{-\alpha}(y)=p_{-\alpha}\left(\frac y{1+axy}\right)\alpha^*(1+axy)p_\alpha\left(\frac x{1+axy}\right).\]
此外 $a$ \Term{invertible}{可逆}，即 $a\in\mathbf G_m(S)$，且 $\alpha\circ\alpha^*=2$。
\end{SGAresult}
\SourcePage{36}
\begin{proof}
A) 考虑由 $(x,y)\mapsto p_\alpha(x)p_{-\alpha}(y)$ 定义的 $\mathbf G_{a,S}^2\to G$，以 $U$ 为 $\Omega$ 的\Term{inverse-image}{逆像}。这是包含 $0\times_S\mathbf G_{a,S}$ 与 $\mathbf G_{a,S}\times_S0$ 的\Term{open-subset}{开集}。存在唯一的 $S$-\Term{scheme}{概形}\Term{morphism}{态射} $A,C:U\to\mathbf G_{a,S}$、$B:U\to T$，满足
\[p_\alpha(u)p_{-\alpha}(v)=p_{-\alpha}(A(u,v))B(u,v)p_\alpha(C(u,v)).\]
立即有
\[A(0,v)=v,\quad A(u,0)=0,\quad C(u,0)=u,\quad C(0,v)=0,\quad B(u,0)=B(0,v)=e.\]
设 $S'$ 为\Term{separated}{分离}的 $S$-\Term{scheme}{概形}，$t\in T(S')$。$\Omega_{S'}$ 在 $\operatorname{int}(t)$ 下\Term{stable}{稳定}，故由1.16末式，$U_{S'}$ 在 $(x,y)\mapsto(\alpha(t)x,\alpha(t)^{-1}y)$ 下\Term{stable}{稳定}，且
\[\begin{aligned}A(\alpha(t)u,\alpha(t)^{-1}v)&=\alpha(t)^{-1}A(u,v),\\C(\alpha(t)u,\alpha(t)^{-1}v)&=\alpha(t)C(u,v),\\B(\alpha(t)u,\alpha(t)^{-1}v)&=B(u,v).\end{aligned}\]
$\alpha$ \Term{faithfully-flat}{忠实平坦}，故对任意 $S'\to S$、$z\in\mathbf G_m(S')$，$U_{S'}$ 在 $(x,y)\mapsto(zx,z^{-1}y)$ 下\Term{stable}{稳定}，并有
\[A(zu,z^{-1}v)=z^{-1}A(u,v),\quad C(zu,z^{-1}v)=zC(u,v),\quad B(zu,z^{-1}v)=B(u,v).\]
先设 $v$ \Term{invertible}{可逆}，取 $z=v$，则 $(u,v)\in U(S')$ 蕴含 $(uv,1)\in U(S')$，并且 $A(uv,1)=v^{-1}A(u,v)$、$B(uv,1)=B(u,v)$。令 $V\subset\mathbf G_{a,S}^2$ 为\Term{open-subset}{开集}\footnote{编者注：加入了条件“$(1,uv)\in U(S')$”。}
\[(u,v)\in V(S')\Longleftrightarrow(u,v),(uv,1),(1,uv)\in U(S').\]
$U$ 包含两坐标轴，故 $V$ 为\Term{zero-section}{零截面}的\Term{neighborhood}{邻域}。以上已说明，\Term{morphism}{态射} $(u,v)\mapsto A(u,v)$ 与 $(u,v)\mapsto vA(uv,1)$，以及 $(u,v)\mapsto B(u,v)$ 与 $(u,v)\mapsto B(uv,1)$ 在 $V\cap(\mathbf G_{a,S}\times_S\mathbf G_{m,S})$ 上相同。后者在 $\mathbf G_{a,S}^2$ 中\Term{schematically-dense}{概形稠密}，故这些\Term{morphism}{态射}在 $V$ 上相同。
\SourcePage{37}
由 $A(0,1)=1$，存在 $\mathbf G_{a,S}$ 中包含\Term{zero-section}{零截面}的\Term{open-subset}{开集} $W_1$，使 $A(x,1)$ 对任意 $W_1$ 的\Term{section}{截面} $x$ \Term{invertible}{可逆}。令 $F(x)=A(x,1)^{-1}$，则若 $(u,v)\in V(S')$\footnote{编者注：此处及下文将 $U(S')$ 改为 $V(S')$。}且 $uv\in W_1(S')$，有 $A(u,v)=vA(uv,1)=vF(uv)^{-1}$。同样处理 $C$，得到包含\Term{zero-section}{零截面}的\Term{open-subset}{开集} $W_2$ 及 $E\in\mathscr O(W_2)^\times$\footnote{编者注：将原书的 $G$ 记为 $E$，因为 $G$ 已用于所考察的 $S$-群。}，使 $C(u,v)=uC(1,uv)=uE(uv)^{-1}$。因此令 $W=W_1\cap W_2$，得到：存在包含\Term{zero-section}{零截面}的\Term{open-subset}{开集} $W\subset\mathbf G_{a,S}$ 及 $S$-\Term{morphism}{态射}
\[F:W\to\mathbf G_{m,S},\ F(0)=1;\quad H:W\to T,\ H(0)=e;\quad E:W\to\mathbf G_{m,S},\ G(0)=1,\]
使当 $(u,v)\in V(S')$、$uv\in W(S')$ 时，
\[p_\alpha(u)p_{-\alpha}(v)=p_{-\alpha}(vF(uv)^{-1})H(uv)p_\alpha(uE(uv)^{-1}).\tag{+}\]
B) 用 $G$ 的结合律写成 $p_\alpha(u)p_{-\alpha}(v)p_{-\alpha}(w)=p_\alpha(u)p_{-\alpha}(v+w)$。存在包含\Term{unit-section}{单位截面}的\Term{open-subset}{开集} $L\subset\mathbf G_{a,S}^3$，使 $(u,v,w)\in L(S')$ \Term{equivalence}{等价}于
\[\begin{gathered}(u,v),(uE(uv)^{-1},w),(u,v+w)\in V(S'),\\uv,\ uwE(uv)^{-1},\ u(v+w)\in W(S').\end{gathered}\]
于是由(+)立即得到
\begin{align*}E(uv+uw)&=E(uwE(uv)^{-1})E(uv),\tag{1}\\H(uv+uw)&=H(uwE(uv)^{-1})H(uv),\tag{2}\\(v+w)F(uv+uw)^{-1}&=\alpha(H(uv)^{-1})wF(uwE(uv)^{-1})^{-1}+vF(uv)^{-1}.\tag{3}\end{align*}
由定义 $(1,0,0)\in L(S)$。考虑 $L\cap(1\times_S\mathbf G_{a,S}\times_S\mathbf G_{a,S})=1\times_SM$；$M$ 是包含 $(0,0)$ 的\Term{open-subset}{开集}，且对 $(v,w)\in M(S')$，$v,wE(v)^{-1},v+w\in W(S')$，并有
\begin{align*}E(v+w)&=E(wE(v)^{-1})E(v),\tag{$1'$}\\H(v+w)&=H(wE(v)^{-1})H(v),\tag{$2'$}\\(v+w)F(v+w)^{-1}&=\alpha(H(v))^{-1}wF(wE(v)^{-1})^{-1}+vF(v)^{-1}.\tag{$3'$}\end{align*}
\SourcePage{38}
最后考虑 $M\to\mathbf G_{a,S}^2$，$(v,w)\mapsto(v,wE(v)^{-1})$。\footnote{编者注：修正了以下内容。}它保持 $(0,0)$，并诱导 $M$ 到包含\Term{zero-section}{零截面}的\Term{open-subset}{开集} $N$ 的\Term{isomorphism}{同构}，逆\Term{isomorphism}{同构}为 $(x,y)\mapsto(x,yE(x))$。\footnote{编者注：即作变量替换 $x=v,y=wE(v)^{-1}$，亦即 $v=x,w=yE(x)$。}故已证明：存在包含\Term{zero-section}{零截面}的\Term{open-subset}{开集} $N\subset\mathbf G_{a,S}^2$，使 $(x,y)\in N(S')$ 时，$x,y,x+yE(x)$\footnote{编者注：将 $yE(x)$ 改为 $x+yE(x)$。} 均在 $W(S')$ 中，且
\begin{align*}E(x+yE(x))&=E(x)E(y),\tag{$1''$}\\H(x+yE(x))&=H(x)H(y),\tag{$2''$}\\(x+yE(x))F(x+yE(x))^{-1}&=xF(x)^{-1}+r(H(x))^{-1}yE(x)F(y)^{-1}.\tag{$3''$}\end{align*}
C) 对左结合律作相同论证，得出：\footnote{编者注：写出 $p_\alpha(t)p_\alpha(u)p_{-\alpha}(v)=p_\alpha(t+u)p_{-\alpha}(v)$ 所给的等式，再令 $v=1,x=u,t=yF(u)$（即 $y=tF(u)^{-1}$）。}存在包含\Term{zero-section}{零截面}的\Term{open-subset}{开集} $N'\subset\mathbf G_{a,S}^2$，使 $(x,y)\in N(S')$ 时，$x,y,x+yF(x)$\footnote{编者注：将 $yF(x)$ 改为 $x+yF(x)$。} 均在 $W(S')$ 中，且
\begin{align*}F(x+yF(x))&=F(x)F(y),\tag{$4''$}\\H(x+yF(x))&=H(x)H(y),\tag{$5''$}\\(x+yF(x))E(x+yF(x))^{-1}&=xE(x)^{-1}+\alpha(H(x))^{-1}yF(x)E(y)^{-1}.\tag{$6''$}\end{align*}
于是须解“函数方程”$(1'')$。
\end{proof}
\begin{SGAresult}{引理}{2.3}
设 $S$ 为\Term{scheme}{概形}，$W$ 为 $\mathbf G_{a,S}$ 中包含\Term{unit-section}{单位截面}的\Term{open-subset}{开集}，$F:W\to\mathbf G_{m,S}$ 为 $S$-\Term{morphism}{态射}，$F(0)=1$。假设存在包含\Term{zero-section}{零截面}的\Term{open-subset}{开集} $N\subset\mathbf G_{a,S}^2$，使 $(x,y)\in N(S')$ 时，$x,y,x+yF(x)$ 均在 $W(S')$ 中，且
\[F(x+yF(x))=F(x)F(y).\tag{$\dagger$}\]
(i) 若 $S$ 是域 $k$ 的谱，则存在 $a\in k$，使 $F(x)=1+ax$。

(ii) 若 $a=F'(0)\in\Gamma(S,\mathscr O_S)$ \Term{invertible}{可逆}，则 $F(x)=1+ax$。
\end{SGAresult}
\begin{proof}
在 $x=0$（相应地 $y=0$）处求导，得
\[F'(y)(1+yF'(0))=F'(0)F(y)\quad((0,y)\in N(S')),\tag{*}\]
相应地 $F'(x)F(x)=F(x)F'(0)$，其中 $(x,0)\in N(S')$。$F$ 取值于 $\mathbf G_m$，故
\[F'(x)=F'(0)\quad((x,0)\in N(S')).\tag{$*'$}\]
\SourcePage{39}
结合第一式，得到
\[F'(0)(1+yF'(0))=F'(0)F(y)\quad((y,0),(0,y)\in N(S')).\]
若 $a=F'(0)$ \Term{invertible}{可逆}，则 $F(y)=1+ay$，对 $W$ 中某个包含\Term{unit-section}{单位截面}的\Term{open-subset}{开集}的\Term{section}{截面} $y$ 成立。此\Term{open-subset}{开集}在 $W$ 中\Term{schematically-dense}{概形稠密}，故(ii)成立。这也证明 $F'(0)\ne0$ 时的(i)。

若 $F'(0)=0$，由$(*')$，在“零附近” $F'(x)=0$，由\Term{schematically-dense}{概形稠密}性得 $F'=0$。若 $k$ \Term{characteristic}{特征}为零，则 $F$ 是导数为零的有理分式，故恒为 $F(0)=1$。

若 $k$ \Term{characteristic}{特征}为 $p$ 而 $F$ 非常数，\footnote{编者注：修正了以下内容。}则存在整数 $n>0$ 及 $F_1\in k(X)$，$F'_1(X)\ne0$，使
\[F(X)=F_1(X^{p^n})=F_1(X)^{p^n}.\]
代回函数方程得
\[F_1(x+yF_1(x)^{p^n})=F_1(x)F_1(y).\tag{$\dagger_1$}\]
在 $x=0$ 处求导得 $F'_1(y)=F'_1(0)F_1(y)$；在 $y=0$ 处求导得 $F'_1(x)F_1(x)^{p^n}=F_1(x)F'_1(0)$。由于 $F'_1(X)$ 在 $k(X)$ 中\Term{invertible}{可逆}，两式给出 $F_1(X)^{p^n}=1$，故 $F_1$ 为常数，与起初假设矛盾。故 $F$ 恒为 $1=F(0)$。
\end{proof}
\begin{proof}[2.2的证明（续）]
D) 设 $S$ 为域的谱。若 $F'(0)=0$，则 $F=1$，由$(5'')$，$H(x+y)=H(x)H(y)$，故 $H$ 延拓为\Term{group-morphism}{群态射} $\mathbf G_{a,S}\to T$（第 XVIII 讲2.3），必恒为 $e$。又由2.3，$E(x)=1+bx$，某个 $b\in k$。于是$(6'')$ 对 $(x,y)\in N(S')$ 给出
\[(x+y)E(x+y)^{-1}=xE(x)+yE(y)^{-1}.\]
再由第 XVIII 讲2.3，$x\mapsto xE(x)^{-1}$ 延拓为 $k$-\Term{group-morphism}{群态射} $\mathbf G_{a,k}\to\mathbf G_{a,k}$，\footnote{编者注：添加了下一句。也可直接计算，上一等式蕴含 $0=xyb(2+(x+y)b)$，从而 $0=b(2+(x+y)b)$，最终 $b=0$。}故 $x/(1+bx)=cx$，某个 $c\in k$，遂有 $b=0$（且 $c=1$）。因此 $F,H,E$ 在\Term{unit-section}{单位截面}\Term{neighborhood}{邻域}恒为 $(1,e,1)$，故处处如此；由(+)可知 $U_\alpha,U_{-\alpha}$ 交换，与(ii)矛盾。
\SourcePage{40}
现在对任意 $S$，已证明 $F'(0)$ 在任何\Term{fiber}{纤维}上均不为零，故\Term{invertible}{可逆}。$E'(0)$ 显然也如此，故由2.3存在 $a,b\in\mathbf G_m(S)$，使
\[F(x)=1+ax,\quad E(x)=1+bx\quad(x\in W(S')).\tag{$\diamond_1$}\]
E) 余下部分很容易。代入$(3'')$得到
\[y\alpha(H(x))(1+ay)=y\bigl(1+ax+ay(1+bx)\bigr)(1+bx).\]
它对 $N$ 的任意\Term{section}{截面}成立。由 $\mathbf G_{a,S}\times_S\mathbf G_{m,S}$ 在 $\mathbf G_{a,S}^2$ 中\Term{schematically-dense}{概形稠密}，有
\[(1+ay)\alpha(H(x))=\bigl(1+ax+ay(1+bx)\bigr)(1+bx).\]
取 $y=0$，得 $\alpha(H(x))=(1+ax)(1+bx)$。代回上式，\footnote{编者注：并利用 $1+bx$ \Term{invertible}{可逆}。}得 $a^2xy=abxy$。由 $\mathbf G_{m,S}$ 在 $\mathbf G_{a,S}$ 中\Term{schematically-dense}{概形稠密}，$a^2=ab$；由于 $a$ \Term{invertible}{可逆}，
\[a=b,\qquad\alpha(H(x))=(1+ax)^2.\tag{$\diamond_2$}\]
$a$ \Term{invertible}{可逆}，故 $x\mapsto1+ax$ 是 $\mathbf G_{a,S}$ 的\Term{automorphism}{自同构}。因此存在包含\Term{section}{截面} $1$ 的\Term{open-subset}{开集} $W'\subset\mathbf G_{a,S}$ 和\Term{morphism}{态射} $P:W'\to T$，使 $P(1+ax)=H(x)$。\footnote{编者注：即作变量替换 $x'=1+ax$，亦即 $x=(x'-1)/a$。}代入$(2')$，对 $(x,y)\in N(S')$ 得
\[P(1+ax+ay)=P\left(\frac{1+ax+ay}{1+ax}\right)P(1+ax).\]
因此在 $\mathbf G_{m,S}$ 的 $1$ 的某\Term{open-neighborhood}{开邻域}内，$P(x)P(y)=P(xy)$。由第 XVIII 讲2.3，存在延拓 $P$ 的\Term{group-morphism}{群态射}
\[\alpha^*:\mathbf G_{m,S}\longrightarrow T.\tag{$\diamond_3$}\]
$\alpha(H(x))=(1+ax)^2$ 在\Term{zero-section}{零截面}附近成立，故 $\alpha(\alpha^*(z))=z^2$ 在\Term{section}{截面} $1$ 附近成立，从而
\[\alpha\circ\alpha^*=2.\tag{$\diamond_4$}\]
\SourcePage{41}
F)\footnote{编者注：略改了以下内容，因为A)中已经引入\Term{open-subset}{开集} $V$。}综合(+)与$(\diamond_1$--$\diamond_4)$，存在 $a\in\mathbf G_m(S)$ 及 $\alpha^*\in\operatorname{Hom}_{S\text{-gr.}}(\mathbf G_{m,S},T)$，$\alpha\circ\alpha^*=2$，使若 $(u,v)\in V(S')$、$uv\in W(S')$，则 $1+auv$ \Term{invertible}{可逆}，且
\[p_\alpha(u)p_{-\alpha}(v)=p_{-\alpha}\left(\frac v{1+auv}\right)\alpha^*(1+auv)p_\alpha\left(\frac u{1+auv}\right).\]
令 $V'=(\mathbf G_{a,S}^2)_f$，$f(u,v)=1+auv$。两端定义 $V'\to G$ 的\Term{morphism}{态射}，在\Term{zero-section}{零截面}\Term{neighborhood}{邻域}相同，故在 $V'$ 相同。因此公式对 $V'$ 任意\Term{section}{截面}成立，特别地 $V'\subset U$，$U$ 如A)所定义。

证明 $U=V'$。回到A)记号，$A:U\to\mathbf G_{a,S}$ 在 $V'$ 上为 $A(u,v)=v(1+auv)^{-1}$。等式 $U=V'$ 是\Term{set-theory}{集合论}问题，可归结于 $S$ 为域 $k$ 的谱；这时就是以下显然事实：由有理分式 $Y/(1+aXY)$ 定义的\Term{rational-map}{有理映射} $\mathbf G_{a,k}^2\dashrightarrow\mathbf G_{a,k}$ 的定义域正是由 $1+aXY$ 定义的\Term{open-subset}{开集}。

G) 现已证明 $a,\alpha^*$ 的存在及两项附\Term{additive}{加性}质。尚须证明唯一性。设 $a',\alpha'^*$ 亦满足要求。对 $u,v\in\mathbf G_a(S')$，立即有
\[1+auv\text{ 可逆}\ \Longrightarrow\ 1+a'uv\text{ 可逆，且 }\frac v{1+auv}=\frac v{1+a'uv}.\]
故对任意 $u\in\mathbf G_a(S')$，$1+au$ \Term{invertible}{可逆}蕴含 $1+au=1+a'u$，立即给出 $a=a'$。再有 $1+au$ \Term{invertible}{可逆}蕴含 $\alpha^*(1+au)=\alpha'^*(1+au)$，故 $\alpha^*=\alpha'^*$。
\end{proof}
\begin{SGAresult}{推论}{2.4}
设 $\exp(Y)t\exp(X)$、$\exp(Y')t'\exp(X')$ 为 $\Omega(S')$ 的两元素。其积在 $\Omega(S')$ 中的充要条件是 $u=1+\langle X,Y'\rangle$ \Term{invertible}{可逆}；此时
\[\begin{split}&\exp(Y)t\exp(X)\exp(Y')t'\exp(X')\\&=\exp(Y+u^{-1}\alpha(t)^{-1}Y')\,tt'\alpha^*(u)\,\exp(u^{-1}\alpha(t')^{-1}X+X').\end{split}\tag{$F'$}\]
\end{SGAresult}
\SourcePage{42}
\textbf{注记 2.5。}定理2.1的(F)也可不借助 $\exp$ 来写。以这些\Term{morphism}{态射}搬运\Term{duality}{对偶} $\mathfrak g^\alpha\otimes\mathfrak g^{-\alpha}\to\mathscr O_S$，得到\Term{vector-bundle}{向量丛}的\Term{canonical}{典范}\Term{pairing}{配对} $U_\alpha\times_SU_{-\alpha}\to\mathbf G_{a,S}$，仍记为 $(x,y)\mapsto\langle x,y\rangle$。于是 $\langle\exp X,\exp Y\rangle=\langle X,Y\rangle$。若 $x\in U_\alpha(S')$、$y\in U_{-\alpha}(S')$ 且 $1+\langle x,y\rangle\in\mathbf G_m(S')$，则
\[x\cdot y=y^{(1+\langle x,y\rangle)^{-1}}\alpha^*(1+\langle x,y\rangle)x^{(1+\langle x,y\rangle)^{-1}}.\tag{F}\]
\begin{SGAresult}{推论}{2.6}
\Term{pairing}{配对} $\mathbf W(\mathfrak g^\alpha)\times_S\mathbf W(\mathfrak g^{-\alpha})\to\mathbf G_{a,S}$ 定义 $\mathbf G_{m,S}$ 下\Term{principal-bundle}{主丛}的\Term{pairing}{配对}
\[\mathbf W(\mathfrak g^\alpha)^\times\times_S\mathbf W(\mathfrak g^{-\alpha})^\times\longrightarrow\mathbf G_{m,S}.\]
\end{SGAresult}
仍将此\Term{pairing}{配对}记为 $(X,Y)\mapsto\langle X,Y\rangle$，或简记为 $XY$。每个 $X\in\Gamma(S,\mathfrak g^\alpha)^\times$ 都有唯一 $X^{-1}\in\Gamma(S,\mathfrak g^{-\alpha})^\times$，使 $XX^{-1}=1$。有 $(zX)^{-1}=z^{-1}X^{-1}$，故所定义的 $s:\mathbf W(\mathfrak g^\alpha)^\times\to\mathbf W(\mathfrak g^{-\alpha})^\times$ 为\Term{scheme}{概形}\Term{isomorphism}{同构}，与\Term{operator}{算子}群上的\Term{isomorphism}{同构} $s:z\mapsto z^{-1}$ 相容。
\begin{SGAresult}{定义}{2.6.1}称 $X$ 与 $s(X)=X^{-1}$ \Term{paired}{相配}。\end{SGAresult}
把2.4用于 $Y=0=X'$、$Y'=aX^{-1}$，$a\in\mathscr O_S(S)$，则 $u=1+a$ 且 $u^{-1}Y'=(1-u^{-1})X^{-1}$，故：
\begin{SGAresult}{推论}{2.7}
设 $X\in\Gamma(S,\mathfrak g^\alpha)^\times$、$u\in\Gamma(S,\mathscr O_S)^\times$，则
\[\alpha^*(u)=\exp((u^{-1}-1)X^{-1})\exp(X)\exp((u-1)X^{-1})\exp(-u^{-1}X).\]
\end{SGAresult}
\begin{SGAresult}{定义}{2.8}称 $\alpha^*$ 为\Term{root}{根} $\alpha$ 对应的\Term{coroot}{余根}。\end{SGAresult}
\textbf{注记 2.9。}若 $(G,T,\alpha)$ 为 $S$-\Term{elementary-system}{初等系统}，$(G,T,-\alpha)$ 亦然，由2.1得到 $\mathfrak g^{-\alpha}$ 与 $\mathfrak g^\alpha$ 的\Term{duality}{对偶}及\Term{coroot}{余根} $(-\alpha)^*$。对(F)取逆，立即有
\[\langle X,Y\rangle=\langle Y,X\rangle,\qquad(-\alpha)^*=-\alpha^*.\]
现考虑 $G$ 的\Term{lie-algebra}{李代数}。\Term{root}{根} $\alpha$ 与\Term{coroot}{余根} $\alpha^*$ 定义\Term{linear-form}{线性形式}
\[\mathscr O_S\xrightarrow{\alpha_*}\mathfrak t\xrightarrow{\alpha}\mathscr O_S.\]
记 $H_\alpha=\alpha_*(1)$。称 $\alpha$ 为对应于 $\alpha$ 的\Term{infinitesimal-root}{无穷小根}，$H_\alpha$ 为相应的\Term{infinitesimal-coroot}{无穷小余根}。
\SourcePage{43}
\begin{SGAresult}{引理}{2.10}
设 $S'\to S$，$X,X'\in\mathbf W(\mathfrak g^\alpha)(S')$、$H\in\mathbf W(\mathfrak t)(S')$、$Y,Y'\in\mathbf W(\mathfrak g^{-\alpha})(S')$、$t\in T(S')$。则
\[\operatorname{Ad}(t)H=H,\quad\operatorname{Ad}(t)X=\alpha(t)X,\quad\operatorname{Ad}(t)Y=\alpha(t)^{-1}Y.\tag{1}\]
\[\begin{gathered}\operatorname{Ad}(\exp X)H=H-\alpha(H)X,\quad\operatorname{Ad}(\exp X)X'=X',\\\operatorname{Ad}(\exp X)Y=Y+\langle X,Y\rangle H_\alpha-\langle X,Y\rangle X.\end{gathered}\tag{2}\]
\[\begin{gathered}\operatorname{Ad}(\exp Y)H=H+\alpha(H)Y,\quad\operatorname{Ad}(\exp Y)Y'=Y',\\\operatorname{Ad}(\exp Y)X=X+\langle X,Y\rangle H_{-\alpha}-\langle X,Y\rangle Y.\end{gathered}\tag{$2'$}\]
\[[H,X]=\alpha(H)X,\quad[H,Y]=-\alpha(H)Y,\quad[X,Y]=\langle X,Y\rangle H_\alpha.\tag{3}\]
\[H_{-\alpha}=-H_\alpha.\tag{4}\]\[\alpha(H_\alpha)=2.\tag{5}\]
\end{SGAresult}
这些公式或者显然，或者直接由2.1的(F)推出。
\begin{SGAresult}{推论}{2.11}
设 $H_\alpha$ 在每个\Term{fiber}{纤维}上非零（特别地，由(5)，$2$ 在 $S$ 上\Term{invertible}{可逆}时如此）。则 $X_\alpha\in\Gamma(S,\mathfrak g^\alpha)^\times$ 与 $X_{-\alpha}\in\Gamma(S,\mathfrak g^{-\alpha})^\times$ \Term{paired}{相配}当且仅当 $[X_\alpha,X_{-\alpha}]=H_\alpha$。
\end{SGAresult}
\textbf{2.12。}设 $(G,T,\alpha)$ 为 $S$-\Term{elementary-system}{初等系统}。由1.19，$\operatorname{Centr}(G)=\operatorname{Ker}(\alpha)$ 为\Term{finite-type}{有限型}\Term{multiplicative-type}{乘法型}群。若 $Q$ 为\Term{center}{中心}的\Term{multiplicative-type}{乘法型}\Term{subgroup}{子群}，则 $G/Q$ 在 $S$ 上\Term{affine}{仿射}（第 IX 讲2.5）、\Term{smooth}{光滑}（第 VI$_B$ 讲9.2），\Term{fiber}{纤维}\Term{connected}{连通}且为\Term{semisimple-rank}{半单秩} $1$ 的\Term{reductive-group}{约化群}（第 XIX 讲1.8）。

令 $G'=G/Q$，它是\Term{semisimple-rank}{半单秩} $1$ 的 $S$-\Term{reductive-group}{约化群}，$T'=T/Q$ 是其\Term{maximal-torus}{极大环面}。\Term{open-subset}{开集} $U_{-\alpha}\cdot T\cdot U_\alpha$ 在 $Q$ 下\Term{stable}{稳定}，其商显然\Term{isomorphism}{同构}于 $U_{-\alpha}\times_S(T/Q)\times_SU_\alpha$。记 $\alpha'$ 为 $\alpha$ 诱导的 $T'$ 的\Term{character}{特征标}，则\Term{canonical}{典范}\Term{morphism}{态射} $G\to G'$ 的\Term{derived-morphism}{导态射}\Term{induction-isomorphism}{诱导同构}
\[\mathfrak g^\alpha\xrightarrow\sim\mathfrak g'^{\alpha'},\qquad\mathfrak g^{-\alpha}\xrightarrow\sim\mathfrak g'^{-\alpha'}.\]
特别地，$\alpha'$ 为 $G'$ 相对于 $T'$ 的\Term{root}{根}。将诱导的 $T/Q\to\mathbf G_{m,S}$ 记为 $\alpha/Q$，得到：
\begin{SGAresult}{引理}{2.13}
若 $Q$ 是 $\operatorname{Ker}(\alpha)$ 的\Term{multiplicative-type}{乘法型}\Term{subgroup}{子群}，则 $(G/Q,T/Q,\alpha/Q)$ 为\Term{elementary-system}{初等系统}。
\end{SGAresult}
\SourcePage{44}
\begin{SGAresult}{引理}{2.14}
在以上条件下，下列图交换：
\[\begin{tikzcd}\mathbf W(\mathfrak g^\alpha)\arrow[r,"\exp"]\arrow[d,"\mathrm{can}"']&G\arrow[d,"\mathrm{can}"]&\mathbf W(\mathfrak g^{-\alpha})\arrow[l,"\exp"']\arrow[d,"\mathrm{can}"]\\\mathbf W(\mathfrak g'^{\alpha'})\arrow[r,"\exp"']&G'&\mathbf W(\mathfrak g'^{-\alpha'})\arrow[l,"\exp"]\end{tikzcd}\]
\[\begin{tikzcd}\mathfrak g^\alpha\otimes\mathfrak g^{-\alpha}\arrow[r,"\sim"]\arrow[d,"\mathrm{can}"']&\mathscr O_S\arrow[d,"\mathrm{id}"]\\\mathfrak g'^{\alpha'}\otimes\mathfrak g'^{-\alpha'}\arrow[r,"\sim"']&\mathscr O_S\end{tikzcd}\qquad\begin{tikzcd}&T\arrow[dr,"\alpha"]\arrow[dd,"\mathrm{can}"]&\\\mathbf G_{m,S}\arrow[ur,"\alpha^*"]\arrow[dr,"\alpha'^*"']&&\mathbf G_{m,S}\\&T'\arrow[ur,"\alpha'"']&\end{tikzcd}\]
\end{SGAresult}
\SGAsection{3}{Weyl 群}
\textbf{记号 3.0。}\footnote{编者注：为便于后续引用添加了编号3.0。}若 $(G,T,\alpha)$ 为 $S$-\Term{elementary-system}{初等系统}，记
\[N=\operatorname{Norm}_G(T),\qquad W=\operatorname{Norm}_G(T)/T\]
（参见第 XIX 讲6.3）；$N$ 为 $G$ 的\Term{closed-subgroup}{闭子群}，在 $S$ 上\Term{smooth}{光滑}。记 $N^\times=N-T$ 为 $N$ 在 $T$ 的补集上诱导的开子\Term{scheme}{概形}。\footnote{编者注：将 $Q$ 改为更具提示性的 $N^\times$。}以 $R$ 为 $\operatorname{Ker}(\alpha)$ 的唯一\Term{maximal-torus}{极大环面}，$T_0$ 为 $\alpha^*:\mathbf G_{m,S}\to T$ 的像，它是 $T$ 的一维\Term{subtorus}{子环面}。$T$ 的乘法诱导 $T_0\times_SR\to T$，此\Term{morphism}{态射}\Term{surjectivity}{满射}（故\Term{faithfully-flat}{忠实平坦}）；只须在\Term{geometric-fibre}{几何纤维}上验证，而这立即来自 $\alpha\circ\alpha^*=2$。
\begin{SGAresult}{定理}{3.1}
(i) $W$ \Term{isomorphism}{同构}于\Term{constant-group}{常值群} $(\mathbf Z/2\mathbf Z)_S$。
\SourcePage{45}
(ii) $N^\times$ 是 $T$ 下\Term{locally-trivial}{局部平凡}的\Term{principal-homogeneous-bundle}{主齐性丛}，左\Term{action}{作用}为 $(t,q)\mapsto tq$（相应地\Term{right-action}{右作用}为 $(q,t)\mapsto qt$）。

(iii) 对 $S'\to S$、$w\in N^\times(S')$、$t\in T(S')$，
\[\operatorname{int}(w)t=t\cdot\alpha^*(\alpha(t)^{-1}).\]
在\Term{resolution}{分解} $T_{S'}=(T_0)_{S'}\cdot R_{S'}$ 中，$\operatorname{int}(w)$ 在 $R_{S'}$ 上诱导恒等，在 $(T_0)_{S'}$ 上诱导对称。有
\[\alpha\circ\operatorname{int}(w)=\alpha^{-1},\qquad\operatorname{int}(w)\circ\alpha^*=(\alpha^*)^{-1}.\]
(iv) 对 $S'\to S$、$X\in\mathbf W(\mathfrak g^\alpha)^\times(S')$，令
\[w_\alpha(X)=\exp(X)\exp(-X^{-1})\exp(X).\]
则 $w_\alpha(X)\in N^\times(S')$，所定义的\Term{morphism}{态射} $w_\alpha:\mathbf W(\mathfrak g^\alpha)^\times\to N^\times$ 满足
\[w_\alpha(zX)=\alpha^*(z)w_\alpha(X)=w_\alpha(X)\alpha^*(z)^{-1}\]
（$z\in\mathbf G_m(S')$）。

(v) 有关系 $w_\alpha(X)w_\alpha(Y)=w_\alpha(-XY^{-1})$。特别地，
\[w_\alpha(X)^2=\alpha^*(-1)\in{}_2T(S)\cap\operatorname{Centr}(G)(S),\qquad w_\alpha(X)^{-1}=w_\alpha(-X)=\alpha^*(-1)w_\alpha(X).\]
(vi) 若同样对 $Y\in\mathbf W(\mathfrak g^{-\alpha})^\times(S')$ 定义 $w_{-\alpha}(Y)=\exp(Y)\exp(-Y^{-1})\exp(Y)$，则除与前面类似的公式外，还有
\[w_{-\alpha}(X^{-1})=w_\alpha(X)^{-1}=w_\alpha(-X),\qquad w_\alpha(X)w_{-\alpha}(Y)=\alpha^*(XY).\]
\end{SGAresult}
\begin{proof}
(i) 已见第 XIX 讲2.4，立即推出 $N^\times$ 对(ii)所定义的\Term{action}{作用}为 $T$ 下\Term{principal-homogeneous-bundle}{主齐性丛}；\Term{locally-trivial}{局部平凡}性\footnote{编者注：相对于 Zariski \Term{topology}{拓扑}。}尤其可由(iv)推出。

证明(iii)。若 $w\in N^\times(S)$，则 $\alpha\circ\operatorname{int}(w)$ 显然是 $G$ 相对于 $T$ 的\Term{root}{根}，故\Term{local-scheme}{局部}等于 $\alpha$ 或 $-\alpha$；每个\Term{fiber}{纤维}上均为 $-\alpha$（Bible，12-05，命题1的推论之证明），故 $\alpha\circ\operatorname{int}(w)=-\alpha$。搬运结构得
\[-\alpha^*=\operatorname{int}(w)^{-1}\circ\alpha^*=\operatorname{int}(w)\circ\alpha^*,\]
因为 $\operatorname{int}(w)^2=\operatorname{int}(w^2)$，且 $w^2$ 为 $T$ 的\Term{section}{截面}。因此 $\operatorname{int}(w)$ 在 $T_0$ 上诱导对称；$R$ 为\Term{central-subgroup}{中心子群}，故在 $R$ 上诱导恒等。(iii)中的公式定义满足相同性质的 $T\to T$ \Term{morphism}{态射}，故与 $\operatorname{int}(w)$ 相同。
\SourcePage{46}
证明(iv)。依次有
\[\begin{split}w_\alpha(X)tw_\alpha(X)^{-1}&=\exp(X)\exp(-X^{-1})\exp(X)t\exp(-X)\exp(X^{-1})\exp(-X)\\&=\exp(X)\exp(-X^{-1})\exp(X-\alpha(t)X)\exp(\alpha(t)^{-1}X^{-1})\exp(-\alpha(t)X)t.\end{split}\]
由(F)，
\[\exp(-X^{-1})\exp((1-\alpha(t))X)=\exp((\alpha(t)^{-1}-1)X)\alpha^*(\alpha(t)^{-1})\exp(-\alpha(t)^{-1}X^{-1}).\]
代回得
\[\operatorname{int}(w_\alpha(X))t=\exp(\alpha(t)^{-1}X)\alpha^*(\alpha(t)^{-1})\exp(-\alpha(t)X)t=\exp(aX)\alpha^*(\alpha(t)^{-1})t,\]
其中 $a=\alpha(t)^{-1}-(\alpha\circ\alpha^*)(\alpha(t)^{-1})\alpha(t)$。由 $\alpha\circ\alpha^*=2$，$a=0$，故 $w_\alpha(X)\in N^\times(S')$。

再证(iv)的第二断言。有\footnote{编者注：第一等式来自1.5(i)，结合 $\alpha\circ\alpha^*=2$ 得到
\[\alpha^*(z)\exp(X)\alpha^*(z)^{-1}=\exp(z^2X),\qquad\alpha^*(z)\exp(X^{-1})\alpha^*(z)^{-1}=\exp(z^{-2}X).\tag{$\dagger$}\]
第三等式来自(F)，第四等式再次来自$(\dagger)$。类似计算还表明 $w_\alpha(X)\alpha^*(z^{-1})=w_\alpha(zX)$。}
\[\begin{aligned}\alpha^*(z)w_\alpha(X)&=\exp(z^2X)\exp(-z^{-2}X^{-1})\exp(z^2X)\alpha^*(z)\\&=\exp(zX)\exp((z^2-z)X)\exp(-z^{-2}X^{-1})\exp(z^2X)\alpha^*(z)\\&=\exp(zX)\exp(-z^{-1}X^{-1})\alpha^*(z)^{-1}\exp((z^3-z^2)X)\exp(z^2X)\alpha^*(z)\\&=\exp(zX)\exp(-z^{-1}X^{-1})\exp(zX)=w_\alpha(zX).\end{aligned}\]
证明(v)。由上一结果，首式立即由第二式推出；证明后者：
\[\begin{aligned}w_\alpha(X)^2&=\exp(X)\exp(-X^{-1})\exp(2X)\exp(-X^{-1})\exp(X)\\&=\exp(X)\exp(-X^{-1})\exp(X^{-1})\alpha^*(-1)\exp(-2X)\exp(X)\\&=\exp(X)\alpha^*(-1)\exp(-X)=\alpha^*(-1).\end{aligned}\]
因为 $\alpha(\alpha^*(-1))=(-1)^2=1$，也证明了 $\alpha^*(-1)\in\operatorname{Centr}(G)(S)$。

最后证明(vi)。第一断言是第二断言的特例。第二式两端给出 $\mathbf W(\mathfrak g^\alpha)^\times\times_S\mathbf W(\mathfrak g^{-\alpha})^\times\to G$ 的\Term{morphism}{态射}，只须在每个\Term{fiber}{纤维}上\Term{nonempty}{非空}的\Term{open-subset}{开集}上证明相同（第 XVIII 讲1.4），故只须在 $1+XY$ \Term{invertible}{可逆}时验证。
\SourcePage{47}
此时依次有
\[\begin{aligned}w_\alpha(X)w_{-\alpha}(Y)&=\exp(X)\exp(-X^{-1})\exp(X)\exp(Y)\exp(-Y^{-1})\exp(Y)\\&=\exp(X)\exp(-X^{-1})\exp\left(\frac Y{1+XY}\right)\alpha^*(1+XY)\exp\left(\frac X{1+XY}\right)\exp(-Y^{-1})\exp(Y)\\&=\exp(X)\exp\left(\frac{-X^{-1}}{1+XY}\right)\alpha^*(1+XY)\exp\left(\frac{-Y^{-1}}{1+XY}\right)\exp(Y)\\&=\exp(-X^{-2}Y^{-1})\alpha^*\left(\frac{XY}{1+XY}\right)\exp(X+Y^{-1})\alpha^*(1+XY)\exp\left(\frac{-Y^{-1}}{1+XY}\right)\exp(Y)\\&=\exp(-X^{-2}Y^{-1})\alpha^*(XY)\exp\left(\frac{Y^{-1}+X}{(1+XY)^2}\right)\exp\left(\frac{-Y^{-1}}{1+XY}\right)\exp(Y)\\&=\alpha^*(XY)\exp(-Y)\exp(Y)=\alpha^*(XY).\end{aligned}\]
\end{proof}
\begin{SGAresult}{推论}{3.2}
设 $n\in\mathbf Z$、$n\ne0$。对任意 $w\in G(S)$，下列条件\Term{equivalence}{等价}：(i) $w\in N^\times(S)$；(ii) $\operatorname{int}(w)\circ n\alpha^*=-n\alpha^*$（回顾 $(n\alpha^*)(z)=\alpha^*(z)^n$）。
\end{SGAresult}
由3.1(iii)，(i)$\Rightarrow$(ii)；反过来可假设 $N^\times$ 有\Term{section}{截面}，遂归结为：
\begin{SGAresult}{引理}{3.3}当 $n\ne0$ 时，$\operatorname{Centr}_G(n\alpha^*)=T$。\end{SGAresult}
\begin{proof}
$n\alpha^*$ 的像 $T_0$ 是 $G$ 的\Term{subtorus}{子环面}，故由第 XIX 讲2.8，$\operatorname{Centr}_G(n\alpha^*)$ 为包含 $T$ 的\Term{reductive-group}{约化}\Term{subgroup}{子群}。在每个\Term{fiber}{纤维}上它异于 $G_s$，故等于 $T_s$（第 XIX 讲1.6.3\footnote{编者注：$\operatorname{Centr}_{G_s}(n\alpha_s^*)\ne G_s$ 蕴含 $\dim\operatorname{Centr}_{G_s}(n\alpha_s^*)-\dim T_s<2$；而据所引之处，此差为偶数。}）；由于二者均为 $G$ 的\Term{smooth}{光滑}\Term{subgroup}{子群}，得所求等式。
\end{proof}
\textbf{注记 3.4。}$w_\alpha$ 的构造及其\Term{normalization}{正规化} $T$ 的事实只依赖(F)。特别地，若 $G$ 为满足2.2条件的 $S$-群，则 $\operatorname{Norm}_G(T)$ 在每个\Term{fiber}{纤维}上异于 $T$。因此若 $G$ \Term{affine}{仿射}、\Term{fiber}{纤维}\Term{connected}{连通}且满足2.2的条件，它就是\Term{semisimple-rank}{半单秩} $1$ 的\Term{reductive-group}{约化群}：它在\Term{unit-section}{单位截面}附近\Term{smooth}{光滑}，故处处\Term{smooth}{光滑}，可用第 XIX 讲1.11的判别法。

\textbf{3.5。}在叙述下一定理前，先作几点说明。照例将 $\mathfrak g^{-\alpha}$ 认作 $(\mathfrak g^\alpha)^{\otimes-1}$，将 $\operatorname{Hom}_{\mathscr O_S}(\mathfrak g^{-\alpha},\mathfrak g^\alpha)$ 认作 $(\mathfrak g^\alpha)^{\otimes2}$，因而
\[\operatorname{Isom}_{\mathscr O_S\text{-mod.}}(\mathbf W(\mathfrak g^{-\alpha}),\mathbf W(\mathfrak g^\alpha))\simeq\mathbf W((\mathfrak g^\alpha)^{\otimes2})^\times.\]
若 $w\in N^\times(S)$，则 $\operatorname{Ad}(w)$ 交换 $\mathfrak g^\alpha,\mathfrak g^{-\alpha}$（3.1(iii)），故定义\Term{isomorphism}{同构} $a_\alpha(w):\mathfrak g^{-\alpha}\xrightarrow\sim\mathfrak g^\alpha$，
\SourcePage{48}
将其认作\Term{section}{截面} $a_\alpha(w)\in\Gamma(S,(\mathfrak g^\alpha)^{\otimes2})^\times$。该构造与\Term{base-change}{基变换}相容，故定义\Term{morphism}{态射}
\[a_\alpha:N^\times\longrightarrow\mathbf W((\mathfrak g^\alpha)^{\otimes2})^\times,\]
使对 $S'\to S$、$w\in N^\times(S')$、$Y\in\Gamma(S',\mathfrak g^{-\alpha})^\times$，$a_\alpha(w)Y=\operatorname{Ad}(w)Y$。
\begin{SGAresult}{定理}{3.6}
(i) 对任意 $S'\to S$、$w\in N^\times(S')$、$Y\in\mathbf W(\mathfrak g^{-\alpha})(S')$，有
\[\operatorname{int}(w)\exp(Y)=\exp(a_\alpha(w)Y).\]
(ii) $a_\alpha(tw)=\alpha(t)a_\alpha(w)$，$a_\alpha(wt)=\alpha(t)^{-1}a_\alpha(w)$。

(iii) 同样定义 $a_{-\alpha}:N^\times\to\mathbf W((\mathfrak g^{-\alpha})^{\otimes2})^\times$，则 $a_{-\alpha}(w)=a_\alpha(w)^{-1}$。\footnote{编者注：即 $a_{-\alpha}(w)$ 与 $a_{-\alpha}(w)$ \Term{paired}{相配}，参见2.6.1。}

(iv) 对任意 $S'\to S$、$X\in\mathbf W(\mathfrak g^\alpha)^\times(S')$，$a_\alpha(w_\alpha(X))=-X^2$。
\end{SGAresult}
\begin{proof}
由1.5对 $\exp$ 的刻画，(i)显然；(ii)、(iii)亦直接。证明(iv)。设 $X\in\Gamma(S',\mathfrak g^\alpha)^\times$、$Z\in\Gamma(S',\mathfrak g^\alpha)$。由定义\footnote{编者注：展开了以下内容的细节。}
\[a_\alpha(w_\alpha(X))^{-1}(Z)=\operatorname{Ad}(w_\alpha(X))(Z)=\operatorname{Ad}(\exp X)\operatorname{Ad}(\exp(-X^{-1}))\operatorname{Ad}(\exp X)(Z).\]
由2.10的$(2')$、(2)及(4)、(5)，并用 $\langle X,X^{-1}\rangle=1$（2.6），右端依次为
\[\begin{aligned}\operatorname{Ad}(\exp X)\operatorname{Ad}(\exp(-X^{-1}))(Z)&=\operatorname{Ad}(\exp X)\bigl(Z+\langle X^{-1},Z\rangle(H_\alpha-X^{-1})\bigr)\\&=Z+\langle X^{-1},Z\rangle(H_\alpha-2X-X^{-1}-H_\alpha+X)\\&=Z-\langle X^{-1},Z\rangle X-\langle X^{-1},Z\rangle X^{-1}.\end{aligned}\]
但 $Z=\langle X^{-1},Z\rangle X$，$\langle X^{-1},Z\rangle X^{-1}=X^{-2}Z$，故 $a_\alpha(w_\alpha(X))^{-1}=-X^{-2}$，得证。
\end{proof}
\begin{SGAresult}{推论}{3.7}
特别地，$\operatorname{int}(w_\alpha(X))\exp(X)=\exp(-X^{-1})$，故由 $w_\alpha$ 的定义，
\[w_\alpha(X)\exp(X)w_\alpha(X)^{-1}=\exp(-X)w_\alpha(X)\exp(-X),\]
即经直接计算，$(w_\alpha(X)\exp(X))^3=e$。
\end{SGAresult}
\begin{SGAresult}{推论}{3.8}
设 $X\in\Gamma(S,\mathfrak g^\alpha)^\times$、$n\in\mathbf Z$、$n\ne0$。则 $w_\alpha(X)$ 为唯一满足下列条件的 $w\in G(S)$：
\SourcePage{49}
(i) $\operatorname{int}(w)\circ n\alpha^*=-n\alpha^*$；(ii) $(w\exp(X))^3=e$。
\end{SGAresult}
\begin{proof}
已知 $w_\alpha(X)$ 满足两条件。反之，由3.2及3.1(ii)，存在 $t\in T(S)$，使 $w=w_\alpha(X)t$。记 $u=\exp(X)$，则
\[wuw^{-1}=w_\alpha(X)t\exp(X)t^{-1}w_\alpha(X)^{-1}=\exp(-\alpha(t)X^{-1}),\]
另一方面
\[\begin{aligned}u^{-1}wu^{-1}&=\exp(-X)w_\alpha(X)t\exp(-X)\\&=\exp(-X)w_\alpha(X)\exp(-X)\exp(X-\alpha(t)X)t\\&=\exp(-X^{-1})\exp(X-\alpha(t)X)t=\exp(-X^{-1})t\exp(H).\end{aligned}\]
由 $(wu)^3=e\Longleftrightarrow wuw^{-1}=u^{-1}wu^{-1}$，比较该元素在 $U_{-\alpha}\cdot T\cdot U_\alpha$ 中的两种\Term{resolution}{分解}，得 $t=e$。
\end{proof}
\textbf{注记 3.9。}本节的若干结果可归结为下列左\Term{principal-homogeneous-bundle}{主齐性丛}图：
\[\begin{tikzcd}\mathbf W(\mathfrak g^\alpha)^\times\arrow[r,"w_\alpha"]&N^\times\arrow[r,"a_\alpha"]&\mathbf W((\mathfrak g^\alpha)^{\otimes2})^\times\\\mathbf G_{m,S}\arrow[r,"\alpha^*"']&T\arrow[r,"\alpha"']&\mathbf G_{m,S}.\end{tikzcd}\]
注意 $a_\alpha$ \Term{faithfully-flat}{忠实平坦}（因为 $\alpha$ 如此），而 $w_\alpha$ 为\Term{monomorphism}{单态射}当且仅当 $\alpha^*$ 为\Term{monomorphism}{单态射}。留给读者写出右\Term{principal-bundle}{主丛}结构的相应图以及\Term{root}{根} $-\alpha$ 的同类图，并研究诸图间的关系。
\begin{SGAresult}{引理}{3.10}
设 $S$ 为\Term{scheme}{概形}，$q>0$ 为整数，$x\mapsto x^q$ 定义 $\mathbf G_{a,S}$ 的\Term{endomorphism}{自同态}；$(G,T,\alpha)$、$(G',T',\alpha')$ 为两个 $S$-\Term{elementary-system}{初等系统}，$f:G\to G'$ 为 $S$-\Term{group-morphism}{群态射}。设
\[h:(\mathfrak g^\alpha)^{\otimes q}\xrightarrow\sim\mathfrak g'^{\alpha'},\qquad h^\vee:(\mathfrak g^{-\alpha})^{\otimes q}\xrightarrow\sim\mathfrak g'^{-\alpha'}\]
互为逆配（contragrédient）的 $\mathscr O_S$-模\Term{isomorphism}{同构}。假设对任意 $S'\to S$、$X\in\mathbf W(\mathfrak g^\alpha)(S')$，有 $f(\exp(X))=\exp(h(X^q))$。则下列条件\Term{equivalence}{等价}：

(i) $f(\alpha^*(z))=\alpha'^*(z)^q$；(ii) $f(w_\alpha(Z))=w_{\alpha'}(h(Z^q))$；(iii) $f(\exp(Y))=\exp(h^\vee(Y^q))$。

每个条件均理解为对所有 $S'\to S$、$z\in\mathbf G_m(S')$、$Z\in\mathbf W(\mathfrak g^\alpha)^\times(S')$、$Y\in\mathbf W(\mathfrak g^{-\alpha})(S')$ 成立。
\end{SGAresult}
\SourcePage{50}
\begin{proof}(i)$\Rightarrow$(ii)由3.8，(ii)$\Rightarrow$(iii)由3.7，(iii)$\Rightarrow$(i)由2.7。\end{proof}
\begin{SGAresult}{命题}{3.11}
设 $S$ 为\Term{scheme}{概形}，$a\in\mathbf Z$，$q>0$ 使 $x\mapsto x^q$ 定义 $\mathbf G_{a,S}$ 的\Term{endomorphism}{自同态}；$(G,T,\alpha)$、$(G',T',\alpha')$ 为两个 $S$-\Term{elementary-system}{初等系统}，$f:G\to G'$ 为 $S$-\Term{group-morphism}{群态射}。下列关于 $f$ 的条件\Term{equivalence}{等价}：

(i) $f|_T$ \Term{resolution}{分解}为使下图交换的 $f_T:T\to T'$：
\[\begin{tikzcd}\mathbf G_{m,S}\arrow[r,"\alpha^*"]\arrow[d,"q"']&T\arrow[r,"\alpha"]\arrow[d,"f_T"]&\mathbf G_{m,S}\arrow[d,"q"]\\\mathbf G_{m,S}\arrow[r,"\alpha'^*"']&T'\arrow[r,"\alpha'"']&\mathbf G_{m,S}.\end{tikzcd}\]
(ii) 存在（唯一的）$\mathscr O_S$-模\Term{isomorphism}{同构} $h:(\mathfrak g^\alpha)^{\otimes q}\to\mathfrak g'^{\alpha'}$，使对任意 $S'\to S$、$X\in\mathbf W(\mathfrak g^\alpha)(S')$、$Y\in\mathbf W(\mathfrak g^{-\alpha})(S')$，
\[f(\exp(X))=\exp(h(X^q)),\qquad f(\exp(Y))=\exp(h^\vee(Y^q)).\]
因而 $f$ 也满足3.10的\Term{equivalence}{等价}条件。
\end{SGAresult}
\begin{proof}
(ii)$\Rightarrow$(i)：由3.10，$f\circ\alpha^*=q\alpha'^*$，故由3.3，$f|_T$ 经 $T'$ \Term{resolution}{分解}。尚须证明 $\alpha'(f(t))=\alpha(t)^q$，这立即来自 $f$ 诱\Term{derived-group}{导群}\Term{morphism}{态射} $T\cdot U_\alpha\to T'\cdot U_{\alpha'}$。

证明(i)$\Rightarrow$(ii)。设 $X\in\Gamma(S,\mathfrak g^\alpha)$、$Y\in\Gamma(S,\mathfrak g^{-\alpha})$，令 $p_+(x)=f(\exp(xX))$、$p_-(x)=f(\exp(yY))$；它们为\Term{group-morphism}{群态射} $p_+,p_-:\mathbf G_{a,S}\to G$。有
\[\begin{aligned}\operatorname{int}(\alpha'^*(z))^q(p_+(x))&=\operatorname{int}(f_T(\alpha^*(z)))(f(\exp(xX)))\\&=f(\operatorname{int}(\alpha^*(z))(\exp(xX)))\\&=f(\exp(z^2xX))=p_+(z^2x).\end{aligned}\]
用1.2（$Q=\alpha'^*(\mathbf G_{m,S})$），存在 $X'\in\Gamma(S,\mathfrak g'^{\alpha'})$，使 $f(\exp(xX))=p_+(x)=\exp(x^qX')$。同样存在 $Y'\in\Gamma(S,\mathfrak g'^{-\alpha'})$，使 $f(\exp(yY))=\exp(y^qY')$。$f$ 为\Term{group-morphism}{群态射}，故保持公式(F)，立即得到 $X^qY^q=(XY)^q=X'Y'$。容易推出 $X^q\mapsto X'$、$Y^q\mapsto Y'$ 定义所述\Term{isomorphism}{同构} $h,h^\vee$。
\end{proof}
\SourcePage{51}
\begin{SGAresult}{命题}{3.12}
设 $(G,T,\alpha)$ 为 $S$-\Term{elementary-system}{初等系统}，$w\in Q(S)$，令 $\Omega_0=\Omega\cap\operatorname{int}(w^{-1})(\Omega)$。定义 $\Omega$ 上的函数
\[d(\exp(Y)t\exp(X))=\alpha(t)^{-1}+XY.\]
则 $\Omega_0=\Omega_d$，且对 $\exp(Y)t\exp(X)\in\Omega_0(S')$，令 $z=d(\exp(Y)t\exp(X))$，有
\[\operatorname{int}(w)(\exp(Y)t\exp(X))=\exp(z^{-1}a_\alpha(w)^{-1}X)\,t\alpha^*(z)\,\exp(z^{-1}a_\alpha(w)Y).\tag{$\star$}\]
此外 $d\circ\operatorname{int}(w)=d^{-1}$。
\end{SGAresult}
\begin{proof}
立即有\footnote{编者注：交换原书中的 $a_\alpha(w)$ 与 $a_\alpha(w)^{-1}$ 作了修正，并展开等式 $d\circ\operatorname{int}(w)=d^{-1}$ 的证明。}
\[\begin{aligned}\operatorname{int}(w)(\exp(Y)t\exp(X))&=\exp(a_\alpha(w)Y)t\alpha^*(\alpha(t)^{-1})\exp(a_\alpha(w)^{-1}X)\\&=\exp(a_\alpha(w)Y)\exp(\alpha(t)a_\alpha(w)^{-1}X)t\alpha^*(\alpha(t)^{-1}).\end{aligned}\]
由2.1，这是 $\Omega$ 的\Term{section}{截面}当且仅当 $1+\alpha(t)XY$ \Term{invertible}{可逆}，故 $\Omega_0=\Omega_d$。再用2.1的(F)，直接计算即得$(\star)$。由此又有
\[(d\circ\operatorname{int}(w))(\exp(Y)t\exp(X))=\alpha(t\alpha^*(z))^{-1}+z^{-2}XY=z^{-2}(\alpha(t)^{-1}+XY)=z^{-1}.\]
\end{proof}
\textbf{注意。}函数 $d$ 与 $w$ 的选择无关。
\SGAsection{4}{同构定理}
\begin{SGAresult}{定理}{4.1}
设 $S$ 为\Term{scheme}{概形}，$q\in\mathbf Z$、$q>0$，$x\mapsto x^q$ 是 $\mathbf G_{a,S}$ 的\Term{endomorphism}{自同态}；$(G,T,\alpha)$、$(G',T',\alpha')$ 为两个 $S$-\Term{elementary-system}{初等系统}。设
\[h:(\mathfrak g^\alpha)^{\otimes q}\to\mathfrak g'^{\alpha'},\qquad h^\vee:(\mathfrak g^{-\alpha})^{\otimes q}\to\mathfrak g'^{-\alpha'}\]
为互为逆配的\Term{isomorphism}{同构}。设 $f_T:T\to T'$ 为使下图交换的 $S$-\Term{group-morphism}{群态射}：
\[\begin{tikzcd}\mathbf G_{m,S}\arrow[r,"q"]\arrow[d,"\alpha^*"']&\mathbf G_{m,S}\arrow[d,"\alpha'^*"]\\T\arrow[r,"f_T"]\arrow[d,"\alpha"']&T'\arrow[d,"\alpha'"]\\\mathbf G_{m,S}\arrow[r,"q"']&\mathbf G_{m,S}.\end{tikzcd}\]
\SourcePage{52}
存在唯一延拓 $f_T$ 的 $S$-\Term{group-morphism}{群态射} $f:G\to G'$，使对所有 $S'\to S$、$X\in\mathbf W(\mathfrak g^\alpha)(S')$，$f(\exp(X))=\exp(h(X^q))$。此外，对 $Y\in\Gamma(S',\mathfrak g^{-\alpha})$、$Z\in\Gamma(S',\mathfrak g^\alpha)^\times$，它还满足
\[f(\exp(Y))=\exp(h^\vee(Y^q)),\qquad f(w_\alpha(Z))=w_\alpha(h(Z^q)).\]
\end{SGAresult}
\begin{proof}
若 $f$ 延拓 $f_T$，则 $f\circ\alpha^*=(\alpha'^*)^q$；若还满足第二条件，则由3.10也满足另两条件。故 $f$ 在 $\Omega$ 上由
\[f(\exp(Y)t\exp(X))=\exp(h^\vee(Y^q))f_T(t)\exp(h(X^q))\]
确定。$\Omega$ 在 $G$ 中\Term{schematically-dense}{概形稠密}，已得唯一性。由第 XVIII 讲2.3，存在性只须证明此公式定义 $\Omega\to G'$ 的“\Term{generic-point}{泛点}乘法相容”\Term{morphism}{态射}。由2.4，这归结于验证 $\alpha'\circ f=\alpha^q$，而它来自 $f$ 延拓 $f_T$。
\end{proof}
\textbf{附释 4.2。}也可如下解释4.1：考虑 $S$-\Term{elementary-system}{初等系统}的\Term{category}{范畴} $\mathscr E$，以及由
\[(\mathbf G_{m,S}\xrightarrow{\alpha^*}T\xrightarrow\alpha\mathbf G_{m,S},\mathscr L)\]
组成的\Term{category}{范畴} $\mathscr D$，其中 $T$ 为\Term{torus}{环面}，$\alpha,\alpha^*$ 为满足 $\alpha\circ\alpha^*=2$ 的\Term{group-morphism}{群态射}，$\mathscr L$ 为\Term{invertible-module}{可逆} $\mathscr O_S$-模（请读者明确所考虑的两\Term{category}{范畴}的\Term{morphism}{态射}）。由
\[(G,T,\alpha)\longmapsto(\mathbf G_{m,S}\xrightarrow{\alpha^*}T\xrightarrow\alpha\mathbf G_{m,S},\mathfrak g^\alpha)\]
定义\Term{functor}{函子} $\mathscr E\to\mathscr D$。上一定理说此\Term{functor}{函子}\Term{fully-faithful}{全忠实}。事实上它是\Term{category-equivalence}{范畴等价}，下一节将说明。目前已有：
\begin{SGAresult}{推论}{4.3}若 $q=1$ 且 $f_T$ 为\Term{isomorphism}{同构}，则 $f$ 为\Term{isomorphism}{同构}。\end{SGAresult}
\begin{SGAresult}{推论}{4.4}
若 $q=1$ 且 $f_T$ \Term{faithfully-flat}{忠实平坦}、核为 $Q$（参见第 IX 讲2.7），则 $f$ \Term{faithfully-flat}{忠实平坦}（\Term{quasi-compact}{拟紧}）、核为 $Q$，故将 $G'$ 认同于 $G/Q$。
\end{SGAresult}
\begin{proof}
若 $f_T$ \Term{faithfully-flat}{忠实平坦}、核为 $Q$，则
\[Q=\operatorname{Ker}(f_T)\subset\operatorname{Ker}(f_T\circ\alpha')=\operatorname{Ker}(\alpha).\]
引入2.13的 $S$-\Term{elementary-system}{初等系统} $(G/Q,T/Q,r/Q)$，由2.14归结于证明 $f/Q$ 诱导 $G/Q\xrightarrow\sim G'$，这立即由4.3推出。
\end{proof}
\SourcePage{53}
\SGAsection{5}{初等系统的例子及应用}
\textbf{5.1。}设 $S$ 为\Term{scheme}{概形}，$\mathscr L$ 为\Term{invertible-module}{可逆} $\mathscr O_S$-模。考虑 $S$ 上的群 $G_{\mathscr L}$，定义为
\[G_{\mathscr L}(S')=\left\{\begin{pmatrix}a&b\\c&d\end{pmatrix}\ \middle|\ a,d\in\mathbf G_a(S'),\ b\in\mathbf W(\mathscr L)(S'),\ c\in\mathbf W(\mathscr L^{-1})(S'),\ ad-bc\in\mathbf G_m(S')\right\},\]
配以通常矩阵乘法。它\Term{local-isomorphism}{局部同构}于 $\operatorname{GL}_{2,S}$，故为在 $S$ 上\Term{affine}{仿射}、\Term{smooth}{光滑}且\Term{fiber}{纤维}\Term{connected}{连通}的 $S$-\Term{group-scheme}{群概形}。

\textbf{注记。}设 $\mathscr L',\mathscr L''$ 为 $S$ 上的两个\Term{invertible}{可逆}\Term{sheaf}{层}，$\mathscr L=\mathscr L'\otimes\mathscr L''^{-1}$。\footnote{编者注：将 $\mathscr L''\otimes\mathscr L'^{-1}$ 改为 $\mathscr L'\otimes\mathscr L''^{-1}$，并展开了下一句。}则有 $S$-群\Term{isomorphism}{同构} $G_{\mathscr L}\xrightarrow\sim\operatorname{GL}(\mathscr L'\oplus\mathscr L'')$，如下定义：若 $x$（相应地 $y$）为 $S$ 的\Term{open-subset}{开集} $V$ 上 $\mathscr L'$（相应地 $\mathscr L''$）的\Term{section}{截面}，则
\[\begin{pmatrix}a&b\\c&d\end{pmatrix}\binom xy=\binom{ax+by}{cx+dy}.\]
\textbf{5.2。}以 $S_{\mathscr L}$ 记 $G_{\mathscr L}$ 中由 $ad-bc=1$ 定义的\Term{closed-subgroup}{闭子群}。它也是在 $S$ 上\Term{affine}{仿射}、\Term{smooth}{光滑}且\Term{fiber}{纤维}\Term{connected}{连通}的 $S$-\Term{group-scheme}{群概形}（由以上\Term{isomorphism}{同构}\Term{isomorphism}{同构}于 $\operatorname{SL}(\mathscr L'\oplus\mathscr L'')$）。同样，考虑 $\mathbf G_{m,S}\to G_{\mathscr L}$，$z\mapsto\left(\begin{smallmatrix}z&0\\0&z\end{smallmatrix}\right)$。这是\Term{center}{中心}\Term{monomorphism}{单态射}；取商得到在 $S$ 上\Term{smooth}{光滑}、\Term{affine}{仿射}且\Term{fiber}{纤维}\Term{connected}{连通}的群 $P_{\mathscr L}$（第 VIII 讲5.7）。从上述注记的\Term{isomorphism}{同构}取商，可见 $P_{\mathscr L}$ 认同于\Term{projective-bundle}{射影丛} $\mathbf P(\mathscr L'\oplus\mathscr L'')$ 的\Term{automorphism-group}{自同构群}（参见 EGA II，4.2.7）。以 $i,p$ 记\Term{canonical}{典范}\Term{morphism}{态射}
\[S_{\mathscr L}\xrightarrow iG_{\mathscr L}\xrightarrow pP_{\mathscr L};\]
$i$ 为\Term{closed-immersion}{闭浸入}，$p$ \Term{faithfully-flat}{忠实平坦}且\Term{affine}{仿射}。

\textbf{5.3。}考虑\Term{group-morphism}{群态射}
\[\begin{aligned}t_G:\mathbf G_{m,S}^2\to G_{\mathscr L},&\quad t_G(z,z')=\begin{pmatrix}z&0\\0&z'\end{pmatrix};\\t_S:\mathbf G_{m,S}\to S_{\mathscr L},&\quad t_S(z)=\begin{pmatrix}z&0\\0&z^{-1}\end{pmatrix};\\t_P:\mathbf G_{m,S}\to P_{\mathscr L},&\quad t_P(z)=p(t_G(z,1)).\end{aligned}\]
它们是群\Term{monomorphism}{单态射}，在各群中定义相对\Term{codimension}{余维}为 $2$ 的（\Term{splitting}{分裂}）\Term{torus}{环面}。对每个 $s\in S$，取 $X\in\Gamma(s,\mathscr L\otimes s)^\times$；
\SourcePage{54}
则 $G_{\mathscr L,s}$ 的\Term{section}{截面} $\left(\begin{smallmatrix}0&X\\X^{-1}&0\end{smallmatrix}\right)$ \Term{normalization}{正规化} $t_G(\mathbf G_{m,s}^2)$ 而不\Term{centralizes}{中心化}它。由第 XIX 讲1.6可知 $G_{\mathscr L}$ 为\Term{semisimple-rank}{半单秩} $1$ 的\Term{reductive-group}{约化群}，\Term{maximal-torus}{极大环面}为 $t_G(\mathbf G_{m,S}^2)$。对 $S_{\mathscr L},P_{\mathscr L}$ 作相同论证，得它们均为\Term{semisimple-rank}{半单秩} $1$ 的\Term{reductive-group}{约化群}，\Term{maximal-torus}{极大环面}分别为 $t_S(\mathbf G_{m,S})$、$t_P(\mathbf G_{m,S})$。

\textbf{5.4。}以通常论证立即确定这些群的\Term{lie-algebra}{李代数}及所选\Term{maximal-torus}{极大环面}的\Term{adjoint-action}{伴随作用}。对 $G_{\mathscr L}$，由第 II 讲4.8直接得到
\[\operatorname{Lie}(G_{\mathscr L}/S)=\left\{\begin{pmatrix}a&b\\c&d\end{pmatrix}\ \middle|\ a,d\text{ 为 }\mathscr O_S\text{ 的截面，}b\text{ 为 }\mathscr L\text{ 的截面，}c\text{ 为 }\mathscr L^{-1}\text{ 的截面}\right\},\]
配以通常\Term{lie-bracket}{括号}，且
\[\operatorname{Ad}(t_G(z,z'))\begin{pmatrix}a&b\\c&d\end{pmatrix}=\begin{pmatrix}a&zz'^{-1}b\\z'z^{-1}c&d\end{pmatrix}.\]
记此\Term{lie-algebra}{李代数}为 $\mathfrak g$，令 $\alpha_G:t_G(\mathbf G_{m,S}^2)\to\mathbf G_{m,S}$ 为 $\alpha_G(t_G(z,z'))=zz'^{-1}$ 所定义的\Term{character}{特征标}。由上式立即看出 $\alpha_G$ 是 $G_{\mathscr L}$ 相对于该\Term{torus}{环面}的\Term{root}{根}，而
\[u:\mathscr L\to\mathfrak g,\ X\mapsto\begin{pmatrix}0&X\\0&0\end{pmatrix},\qquad u_-:\mathscr L^{-1}\to\mathfrak g,\ X\mapsto\begin{pmatrix}0&0\\X&0\end{pmatrix}\]
分别为到 $\mathfrak g^{\alpha_G},\mathfrak g^{-\alpha_G}$ 的\Term{isomorphism}{同构}。因而 $(G,t_G(\mathbf G_{m,S}^2),\alpha_G)$ 为 $S$-\Term{elementary-system}{初等系统}。同样令
\[\alpha_S(t_S(z))=z^2,\qquad\alpha_P(t_P(z))=z,\]
则 $(S_{\mathscr L},t_S(\mathbf G_{m,S}),\alpha_S)$、$(P_{\mathscr L},t_P(\mathbf G_{m,S}),\alpha_P)$ 为\Term{elementary-system}{初等系统}，并可定义 $\mathscr L$（相应地 $\mathscr L^{-1}$）到它们\Term{lie-algebra}{李代数}相应\Term{direct-summand}{直和项}的\Term{isomorphism}{同构}。

\textbf{5.5。}令 $\exp\left(\begin{smallmatrix}0&X\\0&0\end{smallmatrix}\right)=\left(\begin{smallmatrix}1&X\\0&1\end{smallmatrix}\right)$。这定义 $\mathbf W(\mathfrak g^{\alpha_G})\to G_{\mathscr L}$，在\Term{lie-algebra}{李代数}上诱导\Term{canonical}{典范}\Term{morphism}{态射}，故是唯一此类\Term{morphism}{态射}（1.5）。同样令 $\exp\left(\begin{smallmatrix}0&0\\Y&0\end{smallmatrix}\right)=\left(\begin{smallmatrix}1&0\\Y&1\end{smallmatrix}\right)$。显式计算(F)得
\[\left\langle\begin{pmatrix}0&X\\0&0\end{pmatrix},\begin{pmatrix}0&0\\Y&0\end{pmatrix}\right\rangle=XY,\qquad\alpha_G^*(z)=\begin{pmatrix}z&0\\0&1/z\end{pmatrix}=t_G(z,z^{-1}).\]
\SourcePage{55}
\Term{open-subset}{开集} $N^\times=N_G^\times$（在3.1前定义）为\footnote{编者注：展开了以下内容的细节。}
\[N_G^\times(S')=\left\{\begin{pmatrix}0&P\\Q&0\end{pmatrix}\ \middle|\ P\in\mathbf W(\mathfrak g^\alpha)^\times(S'),\ Q\in\mathbf W(\mathfrak g^{-\alpha})^\times(S')\right\}.\]
\Term{morphism}{态射} $w_{\alpha_G}$（3.1(iv)）对 $X\in\mathbf W(\mathfrak g^\alpha)^\times(S')$ 为
\[w_{\alpha_G}(X)=\begin{pmatrix}0&X\\-X^{-1}&0\end{pmatrix}.\]
\Term{morphism}{态射} $a_{\alpha_G}$（3.5）为：若 $w=\left(\begin{smallmatrix}0&P\\Q&0\end{smallmatrix}\right)\in N_G^\times(S')$，则
\[a_{\alpha_G}(w)=PQ^{-1}\in\mathbf W((\mathfrak g^\alpha)^{\otimes2})^\times(S'),\]
即对任意 $Y\in\mathbf W(\mathfrak g^{-\alpha})^\times(S')$，$a_{\alpha_G}(w)(Y)=PQ^{-1}Y\in\mathbf W(\mathfrak g^\alpha)^\times(S')$。

\textbf{5.6。}留给读者在 $S_{\mathscr L},P_{\mathscr L}$ 中作同样计算。得到相同的\Term{duality}{对偶}公式及\Term{coroot}{余根}
\[\alpha_S^*(z)=t_S(z),\qquad\alpha_P^*(z)=t_P(z^2).\]
记 $p:G_{\mathscr L}\to P_{\mathscr L}$ 在 $t_S(\mathbf G_{m,S})$ 上诱导的\Term{morphism}{态射}为 $p_T$，即 $p_T(t_S(z))=t_P(z^2)$。于是有\Term{commutative-diagram}{交换图}：\footnote{编者注：其中 $t_S,t_P$ 为\Term{isomorphism}{同构}。}
\[\begin{tikzcd}[column sep=small]& &\mathbf G_{m,S}\arrow[dll,"\mathrm{id}"']\arrow[dl,"\alpha_S^*"']\arrow[dr,"\alpha_P^*"]\arrow[drr,"2"]&&\\\mathbf G_{m,S}\arrow[r,"t_S"]\arrow[drr,"2"']&t_S(\mathbf G_{m,S})\arrow[rr,"p_T"]\arrow[dr,"\alpha_S"']&&t_P(\mathbf G_{m,S})\arrow[dl,"\alpha_P"]&\mathbf G_{m,S}\arrow[l,"t_P"']\arrow[dll,"\mathrm{id}"]\\&&\mathbf G_{m,S}&&\end{tikzcd}\]
中间部分正是4.1中关于\Term{canonical}{典范}\Term{morphism}{态射} $p\circ i:S_{\mathscr L}\to P_{\mathscr L}$ 的\Term{commutative-diagram}{交换图}，\footnote{编者注：取 $q=1$。}它诱导上述 $S$-\Term{elementary-system}{初等系统}的\Term{morphism}{态射}。
\SourcePage{56}
\textbf{5.7。}现设 $(G,T,\alpha)$ 为任意 $S$-\Term{elementary-system}{初等系统}。考虑\Term{commutative-diagram}{交换图}
\[\begin{tikzcd}&\mathbf G_{m,S}\arrow[dl,"\mathrm{id}"']\arrow[d,"\alpha^*"]\arrow[dr,"2"]&\\\mathbf G_{m,S}\arrow[r,"\alpha^*"]\arrow[dr,"2"']&T\arrow[r,"\alpha"]\arrow[d,"\alpha"]&\mathbf G_{m,S}\arrow[dl,"\mathrm{id}"]\\&\mathbf G_{m,S}.&\end{tikzcd}\]
结合以上两图，得到\Term{commutative-diagram}{交换图}
\[\begin{tikzcd}&\mathbf G_{m,S}\arrow[dl,"\alpha_S^*"']\arrow[d,"\alpha^*"]\arrow[dr,"\alpha_P^*"]&\\t_S(\mathbf G_{m,S})\arrow[r,"\alpha^*\circ t_S^{-1}"]\arrow[dr,"\alpha_S"']&T\arrow[r,"t_P\circ\alpha"]\arrow[d,"\alpha"]&t_P(\mathbf G_{m,S})\arrow[dl,"\alpha_P"]\\&\mathbf G_{m,S}.&\end{tikzcd}\]
运用4.1得到：
\begin{SGAresult}{命题}{5.8}
设 $S$ 为\Term{scheme}{概形}，$(G,T,\alpha)$ 为 $S$-\Term{elementary-system}{初等系统}。令 $\mathscr L=\mathfrak g^\alpha$（故 $\mathscr L^{-1}=\mathfrak g^{-\alpha}$）。

(i) 存在唯一\Term{group-morphism}{群态射} $f:S_{\mathscr L}\to G$，满足以下\Term{equivalence}{等价}条件：
\[\begin{aligned}\text{(a)}\quad&f\begin{pmatrix}z&0\\0&z^{-1}\end{pmatrix}=\alpha^*(z),&f\begin{pmatrix}1&X\\0&1\end{pmatrix}&=\exp(X);\\\text{(b)}\quad&f\begin{pmatrix}1&X\\0&1\end{pmatrix}=\exp(X),&f\begin{pmatrix}1&0\\Y&1\end{pmatrix}&=\exp(Y);\\\text{(c)}\quad&f\begin{pmatrix}1&X\\0&1\end{pmatrix}=\exp(X),&f\begin{pmatrix}0&X\\-X^{-1}&0\end{pmatrix}&=w_\alpha(X).\end{aligned}\]
(ii) 存在唯一\Term{group-morphism}{群态射} $g:G\to P_{\mathscr L}$，满足
\[g(t)=\begin{pmatrix}\alpha(t)&0\\0&1\end{pmatrix},\qquad g(\exp(X))=p\begin{pmatrix}1&X\\0&1\end{pmatrix}.\]
此外
\[g(\exp(Y))=p\begin{pmatrix}1&0\\Y&1\end{pmatrix},\qquad g(w_\alpha(X))=p\begin{pmatrix}0&X\\-X^{-1}&0\end{pmatrix}.\]
\SourcePage{57}
$g$ \Term{faithfully-flat}{忠实平坦}、\Term{quasi-compact}{拟紧}，核为 $\operatorname{Ker}(\alpha)=\operatorname{Centr}(G)$，且 $g\circ f$ 是\Term{canonical}{典范}\Term{morphism}{态射} $S_{\mathscr L}\to P_{\mathscr L}$。
\end{SGAresult}
注意(i)中的(b)给出了 $\mathfrak g^\alpha,\mathfrak g^{-\alpha}$ 间\Term{duality}{对偶}的显式描述。
\begin{SGAresult}{推论}{5.9}
设 $(G,T,\alpha)$ 为 $S$-\Term{elementary-system}{初等系统}。\Term{subgroup}{子群} $T\cdot U_\alpha,T\cdot U_{-\alpha},U_\alpha,U_{-\alpha}$ 均为闭的。
\end{SGAresult}
\begin{proof}
$U_\alpha$ 是 $T\cdot U_\alpha$ 的闭子\Term{group-scheme}{群概形}，故只须验证后者。由 Noether 定理（第 IV 讲5.3.1及6.4.1），只须证明 $(T\cdot U_\alpha)/\operatorname{Ker}(\alpha)$ 为 $G/\operatorname{Ker}(\alpha)$ 的\Term{closed-subgroup}{闭子群}。由5.8，归结于证明 $P_{\mathscr L}$（或再次运用 Noether 定理，\Term{equivalence}{等价}地 $G_{\mathscr L}$）中由 $c=0$ 定义的\Term{subgroup}{子群}为闭的，这是显然的。因此1.5(i)中的 $\exp$ 为\Term{closed-immersion}{闭浸入}。
\end{proof}
\textbf{注意。}此推论也来自 $T\cdot U_\alpha,T\cdot U_{-\alpha}$ 为 $G$ 的“Borel\Term{subgroup}{子群}”（第 XII 讲7.10）。

\textbf{5.10。}设 $\mathscr L$ 为\Term{invertible-module}{可逆} $\mathscr O_S$-模，并有群图\footnote{编者注：其中 $T$ 为\Term{torus}{环面}。}
\[\mathbf G_{m,S}\xrightarrow{\alpha^*}T\xrightarrow\alpha\mathbf G_{m,S},\qquad\alpha\circ\alpha^*=2.\]
以 $R$ 为 $\operatorname{Ker}(\alpha)$ 的\Term{maximal-torus}{极大环面}，$K=\alpha^{*-1}(R)$。则 $K$ 为 $\mathbf G_{m,S}$ 的\Term{multiplicative-type}{乘法型}\Term{subgroup}{子群}；由 $\alpha\circ\alpha^*=2$，它甚至是 $\boldsymbol\mu_{2,S}$ 的\Term{subgroup}{子群}。特别地
\[K\longrightarrow S_{\mathscr L},\quad z\longmapsto\begin{pmatrix}z&0\\0&z^{-1}\end{pmatrix}\]
为\Term{center}{中心}\Term{morphism}{态射}。故有\Term{center}{中心}群\Term{monomorphism}{单态射}
\[K\longrightarrow R\times S_{\mathscr L},\quad z\longmapsto\left(\alpha^*(z),\begin{pmatrix}z&0\\0&z^{-1}\end{pmatrix}\right).\]
考虑\Term{quotient-group}{商群} $G=(R\times S_{\mathscr L})/K$。它在 $S$ 上\Term{affine}{仿射}、\Term{smooth}{光滑}且\Term{fiber}{纤维}\Term{connected}{连通}。显然序列
\[1\longrightarrow K\longrightarrow R\times t_S(\mathbf G_{m,S})\xrightarrow uT\longrightarrow1,\qquad u(x,t_S(z))=x\alpha^*(z),\]
\Term{exact}{正合}。因此 $R\times t_S(\mathbf G_{m,S})$ 在 $G$ 中的像是\Term{isomorphism}{同构}于 $T$ 的\Term{torus}{环面} $T'$。不难证明，若 $\alpha'$ 为利用上述\Term{isomorphism}{同构}从 $\alpha$ \Term{derive}{导出}的 $T'$ 的\Term{character}{特征标}，则 $(G,T',\alpha')$ 为 $S$-\Term{elementary-system}{初等系统}，$\mathfrak g^{\alpha'}$ \Term{isomorphism}{同构}于 $\mathscr L$，且 $\alpha'^*$ 经由 $T\xrightarrow\sim T'$ 从 $\alpha^*$ 得到。因此构造了 $S$-\Term{elementary-system}{初等系统} $(G,T',\alpha')$，使相应\Term{object}{对象}
\SourcePage{58}
$(\mathbf G_{m,S}\xrightarrow{\alpha'^*}T'\xrightarrow{\alpha'}\mathbf G_{m,S},\mathfrak g^{\alpha'})$ 在4.2定义的\Term{category}{范畴} $\mathscr D$ 中\Term{isomorphism}{同构}于 $(\mathbf G_{m,S}\xrightarrow{\alpha^*}T\xrightarrow\alpha\mathbf G_{m,S},\mathscr L)$。因此已证明：
\begin{SGAresult}{定理}{5.11}
在4.2的记号下，\Term{functor}{函子}
\[(G,T,\alpha)\longmapsto(\mathbf G_{m,S}\xrightarrow{\alpha^*}T\xrightarrow\alpha\mathbf G_{m,S},\mathfrak g^\alpha)\]
是 $\mathscr E$ 与 $\mathscr D$ 之间的\Term{category-equivalence}{范畴等价}。
\end{SGAresult}
\SGAsection{6}{初等系统的生成元与关系}
\textbf{6.1。}设 $S$ 为\Term{scheme}{概形}，$(G,T,\alpha)$ 为 $S$-\Term{elementary-system}{初等系统}。取 $X\in\mathbf W(\mathfrak g^\alpha)^\times(S)$，令 $u=\exp(X)$。由3.8，$w=w_\alpha(X)$ 特别满足 $(wu)^3=e$。\footnote{编者注：添加了下一句。}以 $s_\alpha$ 记 $\operatorname{int}(w)$ 诱导的 $T$ 的\Term{automorphism}{自同构}；由3.1(iii)，对任意 $S'\to S$、$t\in T(S')$，
\[s_\alpha(t)=\operatorname{int}(w)t=t\cdot\alpha^*(\alpha(t)^{-1}).\]
\begin{SGAresult}{定理}{6.2}
设 $H$ 为(fppf)的 $S$-群\Term{sheaf}{层}，$f_T:T\to H$、$f_\alpha:U_\alpha\to H$ 为\Term{group-morphism}{群态射}，$h\in H(S)$ 为\Term{section}{截面}。存在延拓 $f_T,f_\alpha$ 且满足 $f(w)=h$ 的\Term{group-morphism}{群态射}（必唯一）$f:G\to H$ 的充要条件是：

(i) 对任意 $S'\to S$、$t\in T(S')$、$x\in U_\alpha(S')$，
\[f_T(t)f_\alpha(x)f_T(t)^{-1}=f_\alpha(txt^{-1})=f_\alpha(x^{\alpha(t)}).\tag{1}\]
即 $f_T,f_\alpha$ 可延拓为\Term{semidirect-product}{半直积} $T\cdot U_\alpha$ 到 $H$ 的\Term{group-morphism}{群态射}。

(ii) 对任意 $S'\to S$、$t\in T(S')$，
\[hf_T(t)h^{-1}=f_T(s_\alpha(t))=f_T(t\cdot\alpha^*(\alpha(t)^{-1})).\tag{2}\]
(iii) 在 $H(S)$ 中有两个关系
\begin{align}
h^2&=f_T(\alpha^*(-1)),\tag{3}\\
(hf_\alpha(u))^3&=e.\tag{4}
\end{align}
\end{SGAresult}
\SourcePage{59}
\begin{proof}
用加法记 $U_\alpha,U_{-\alpha}$，以乘法记其向量结构。若 $f$ 满足定理条件，则对任意 $y\in U_{-\alpha}(S')$，必有
\[f(y)=f(w^{-1}wyw^{-1}w)=hf_\alpha(w^{-1}yw)h^{-1}.\]
故定义\Term{group-morphism}{群态射} $f_{-\alpha}:U_{-\alpha}\to H$ 为
\[f_{-\alpha}(y)=hf_\alpha(w^{-1}yw)h^{-1}.\tag{$*_1$}\]
另一方面，$f$ 在\Term{big-cell}{大胞腔} $\Omega$ 上由 $f(ytx)=f_{-\alpha}(y)f_T(t)f_\alpha(x)$ 确定。这证明唯一性。定理条件显然必要，现证充分性。

由(4)，$hf_\alpha(u)h^{-1}h^2=f_\alpha(-u)h^{-1}f_\alpha(-u)$。由(3)及(1)，$h^2=h^{-2}$ 与 $f_\alpha(-u)$ 交换，故
\[hf_\alpha(u)h^{-1}=f_\alpha(-u)hf_\alpha(-u).\]
由定义 $hf_\alpha(u)h^{-1}=f_{-\alpha}(wuw^{-1})$；由3.7，因 $u=\exp(X),w=w_\alpha(X)$，有
\[wuw^{-1}=-\widetilde u,\tag{$*_2$}\]
其中 $\widetilde u$ 为与 $u$ \Term{paired}{相配}的元素。因此
\[f_{-\alpha}(-\widetilde u)=f_\alpha(-u)hf_\alpha(-u).\tag{$*_3$}\]
现取任意 $S'\to S$ 上 $T$ 的\Term{section}{截面} $t$，以 $\operatorname{int}(f_T(t))$ \Term{action}{作用}于此式。左端得到\footnote{编者注：修正了以下内容。}
\[\begin{aligned}f_T(t)f_{-\alpha}(-\widetilde u)f_T(t)^{-1}&=f_T(t)hf_\alpha(u)h^{-1}f_T(t)^{-1}\\&=h(h^{-1}f_T(t)h)f_\alpha(u)(h^{-1}f_T(t)^{-1}h)h^{-1}\\&=hf_T(s_\alpha(t))f_\alpha(u)f_T(s_\alpha(t))^{-1}h^{-1}\\&=hf_\alpha(\alpha(s_\alpha(t))u)h^{-1}\end{aligned}\]
（由(2)、(1)）；因 $s_\alpha(t)=t\alpha^*(\alpha(t)^{-1})$、$\alpha\circ\alpha^*=2$，它等于 $hf_\alpha(\alpha(t)^{-1}u)h^{-1}$。最后由$(*_1)$、$(*_2)$，
\[hf_\alpha(\alpha(t)^{-1}u)h^{-1}=f_{-\alpha}(\alpha(t)^{-1}wuw^{-1})=f_{-\alpha}(-\alpha(t)^{-1}\widetilde u).\]
$(*_3)$ 的右端给出
\[f_\alpha(-\alpha(t)u)\,f_T(t)hf_T(t)^{-1}h^{-1}\,h\,f_\alpha(-\alpha(t)u).\]
由 $hf_T(t)^{-1}h^{-1}=f_T(s_\alpha(t^{-1}))=f_T(t\cdot\alpha^*(\alpha(t)))$，它等于
\[f_\alpha(-\alpha(t)u)f_T(\alpha^*(\alpha(t)))hf_\alpha(-\alpha(t)u).\]
\SourcePage{60}
比较所得两式，得
\[f_{-\alpha}(-\alpha(t)^{-1}\widetilde u)=f_\alpha(-\alpha(t)u)f_T(\alpha^*(\alpha(t)))hf_\alpha(-\alpha(t)u).\]
$\alpha:T\to\mathbf G_{m,S}$ \Term{faithfully-flat}{忠实平坦}，$H$ 为\Term{separated}{分离}预\Term{sheaf}{层}，故可用 $\mathbf G_{m,S}$ 的任意\Term{section}{截面}替换 $-\alpha(t)^{-1}$，得到：
\end{proof}
\begin{SGAresult}{引理}{6.2.1}
对任意 $S'\to S$、$z\in\mathbf G_m(S')$，
\[f_{-\alpha}(z\widetilde u)=f_\alpha(z^{-1}u)f_T(\alpha^*(-z^{-1}))hf_\alpha(z^{-1}u).\]
\end{SGAresult}
\begin{proof}[6.2的证明（续）]
现设 $x,y\in\mathbf G_a(S')$，$y$ 与 $1+xy$ \Term{invertible}{可逆}。先将引理用于 $z=y$，得\footnote{编者注：展开了以下计算。}
\[f_\alpha(xu)f_{-\alpha}(y\widetilde u)=f_\alpha((x+y^{-1})u)f_T(\alpha(-y^{-1}))hf_\alpha(y^{-1}u).\]
有 $x+y^{-1}=y^{-1}(1+xy)$，再将引理用于 $z=y/(1+xy)$，得
\[f_\alpha((x+y^{-1})u)=f_{-\alpha}\left(\frac y{1+xy}\widetilde u\right)f_\alpha(-(x+y^{-1})u)h^{-1}f_T\left(\alpha^*\left(\frac{-y}{1+xy}\right)\right).\]
代入前式，得到
\[f_\alpha(xu)f_{-\alpha}(y\widetilde u)=f_{-\alpha}\left(\frac y{1+xy}u\right)f_\alpha(-(x+y^{-1})u)h^{-1}f_T(\alpha^*((1+xy)^{-1}))hf_\alpha(y^{-1}u).\]
由(2)有 $h^{-1}f_T(t)h=f_T(s_\alpha(t))$（注意 $s_\alpha^2=\mathrm{id}$），且 $s_\alpha\circ\alpha^*=-\alpha^*$（参见6.2.1），故上式等于
\[f_{-\alpha}\left(\frac y{1+xy}u\right)f_\alpha(-(x+y^{-1})u)f_T(\alpha^*(1+xy))f_\alpha(y^{-1}u).\]
最后，对任意 $x'\in U_\alpha(S')$、$z\in\mathbf G_m(S')$，有 $f_\alpha(x')f_T(\alpha^*(z))=f_T(\alpha^*(z))f_\alpha(z^{-2}x')$，故
\[\begin{aligned}f_\alpha(xu)f_{-\alpha}(y\widetilde u)&=f_{-\alpha}\left(\frac y{1+xy}\widetilde u\right)f_T(\alpha^*(1+xy))f_\alpha\left(\left(\frac{-y^{-1}(1+xy)}{(1+xy)^2}+y^{-1}\right)u\right)\\&=f_{-\alpha}\left(\frac y{1+xy}\widetilde u\right)f_T(\alpha^*(1+xy))f_\alpha\left(\frac x{1+xy}u\right).\end{aligned}\]
于是证明了：
\end{proof}
\begin{SGAresult}{引理}{6.2.2}
设 $S'\to S$。若 $a\in U_\alpha(S')$、$b\in U_{-\alpha}^\times(S')$ 且 $1+ab\in\mathbf G_m(S')$，则
\[f_\alpha(a)f_{-\alpha}(b)=f_{-\alpha}\left(\frac b{1+ab}\right)f_T(\alpha^*(1+ab))f_\alpha\left(\frac a{1+ab}\right).\]
\end{SGAresult}
\SourcePage{61}
\begin{proof}[6.2的证明（续）]
由\Term{schematically-dense}{概形稠密}性，此式在 $b\in U_{-\alpha}(S')$ 时仍成立，只须 $1+ab$ \Term{invertible}{可逆}。考虑由 $f(ytx)=f_{-\alpha}(y)f_T(t)f_\alpha(x)$ 定义的\Term{morphism}{态射} $f:\Omega\to H$。由6.2.2、6.2(i)及2.4的$(F')$，立即可知若 $g,g'\in\Omega(S')$ 且 $gg'\in\Omega(S')$，则 $f(gg')=f(g)f(g')$。由第 XVIII 讲2.3(iii)及2.4，\footnote{编者注：注意 $G$ 的每个\Term{geometric-fibre}{几何纤维}均\Term{connected}{连通}，例如由1.1。}存在延拓 $f$ 的\Term{group-morphism}{群态射} $G\to H$，仍记为 $f$。它正是所求：只须证明 $f(w_\alpha)=h$。有 $w_\alpha=u\cdot(-\widetilde u)\cdot u$，故由$(*_3)$，\footnote{编者注：用$(*_3)$简化了原书。}
\[f(w_\alpha)=f_\alpha(u)f_{-\alpha}(-\widetilde u)f_\alpha(u)=h.\]
\end{proof}
\textbf{注记 6.3。}将在第 XXIII 讲3.5补充这些结果。
\SourcePage{62}
% 原 PDF 第62页为空白页。
